% concatenated from main.tex
\documentclass{article}

\usepackage[utf8]{inputenc}
\usepackage[T1]{fontenc}
\usepackage[a4paper, portrait, margin=1in]{geometry}
\usepackage{amsmath,amsfonts,amssymb,amsthm}
\usepackage{booktabs,enumitem,float}
\usepackage{graphicx}
\usepackage{mathpazo}
\usepackage{xcolor}
\usepackage[round]{natbib}
\usepackage{hyperref}
\hypersetup{colorlinks=true, linkcolor=blue, citecolor=teal}
\usepackage[nameinlink, capitalize, noabbrev]{cleveref}
\usepackage[ruled,vlined,linesnumbered]{algorithm2e}
\usepackage{authblk}

\newtheorem{theorem}{Theorem}[section]
\newtheorem{definition}{Definition}[section]
\newtheorem{proposition}{Proposition}[section]
\newtheorem{assumption}[theorem]{Assumption}
\newtheorem{remark}[theorem]{Remark}

\DeclareMathOperator{\Law}{Law}
\DeclareMathOperator{\Cov}{Cov}
\newcommand{\argmin}{\operatorname{arg\,min}}

\crefname{section}{Section}{Sections}
\Crefname{section}{Section}{Sections}
\crefname{subsection}{Section}{Sections}
\Crefname{subsection}{Section}{Sections}
\crefname{appendix}{Appendix}{Appendices}
\Crefname{appendix}{Appendix}{Appendices}
\crefname{algocf}{Algorithm}{Algorithms}
\crefname{assumption}{Assumption}{Assumptions}
\Crefname{assumption}{Assumption}{Assumptions}
\crefname{remark}{Remark}{Remarks}
\Crefname{remark}{Remark}{Remarks}

\newcommand{\cF}{\mathcal{F}}
\newcommand{\cP}{\mathcal{P}}
\newcommand{\cU}{\mathcal{U}}
\newcommand{\bs}{\mathbf{s}}
\newcommand{\bv}{\mathbf{v}}
\newcommand{\Embb}{\mathbb{E}}
\newcommand{\edisp}{\varepsilon_{\mathrm{disp}}}
\newcommand{\emart}{\varepsilon_{\mathrm{mart}}}

\renewcommand\Authsep{,~\,}
\renewcommand\Authands{,~\,}

\title{Global Multi-Maturity SPX--VIX Calibration Beyond Markovian Stitching}
\renewcommand{\thefootnote}{\fnsymbol{footnote}}
\author[1]{Atithi Acharya\thanks{\texttt{atithi.acharya@jpmchase.com}}}
\author[1]{Yue Sun\thanks{\texttt{yue.sun@jpmchase.com}}}
\author[1]{Brandon Augustino}
\author[1]{Shouvanik Chakrabarti}
\author[1]{Shree Hari Sureshbabu}
\author[2]{Charlie Che\thanks{\texttt{charlie.che@jpmchase.com}}}
\affil[1]{Global Technology Applied Research, JPMorganChase, New York, NY 10001, USA}
\affil[2]{Quantitative Trading \& Research, JPMorganChase, New York, NY 10179, USA}
\date{2026}

\begin{document}
\maketitle
\renewcommand{\thefootnote}{\arabic{footnote}}\setcounter{footnote}{0}

\begin{abstract}
We develop a global framework for joint S\&P~500 (SPX)--VIX smile calibration across multiple maturities without the conditional-independence restriction induced by Markovian stitching.
Exact local and global feasibility are equivalent: every globally feasible law has a block-preserving SPX-Markovization that leaves each monthly $(S_i,V_i,S_{i+1})$ law unchanged.
Nevertheless, stitched laws can form a strict subset of globally feasible path laws because Markovization discards dependence on earlier history beyond the current SPX level.
Adjacent smiles therefore cannot identify this dependence, and laws with identical monthly calibrations can price multi-period claims differently.
Under the standard Markov reference, relative entropy selects the stitched minimum-information completion; non-Markov dependence requires cross-period information, an appropriate objective, or a history-dependent prior.
For finite discretizations, we introduce an augmented-Bregman mirror-descent scheme.
It preserves the fit to observable quote moments while controlling martingale and dispersion residuals.
In a controlled infeasible affine system, this split keeps prescribed marginals about $25$ times tighter than cyclic row projection by exposing the discrepancy in the conditional rows.
An exact finite-state example verifies block preservation and exhibits material cross-period price changes after Markovization.
On smoothed SPX and VIX surfaces, numerical calculations illustrate a finite-budget penalty path: the worst fitted-smile error remains below $0.70$ volatility points across the reported sweep while the bulk conditional diagnostics improve substantially.
\end{abstract}

\section{Introduction}\label{sec:intro}

The joint calibration of option-pricing models to S\&P~500 (SPX) and VIX options is a central problem in quantitative finance, driven by the use of volatility derivatives for risk management and hedging.
Accurate joint calibration to both SPX and VIX derivatives is essential for consistent pricing and for the avoidance of arbitrage between desks or across products~\citep{guyon2020joint,cuchiero2025joint}.
The problem is hard because the two markets have conflicting characteristics:
SPX options exhibit a pronounced negative skew, especially at short maturities, while VIX options often display lower implied volatilities at the same maturities, and classical stochastic-volatility models cannot reconcile the two~\citep{carr2006tale,jacquier2018vix}.
A natural nonparametric route is to view joint SPX--VIX calibration through the lens of dispersion-constrained \emph{martingale optimal transport} (MOT)~\citep{Beiglbck2013,guyon2020joint}.
This perspective gives a clean solution in the two-maturity setting, SPX at $T_1$ and $T_2=T_1+30$ days, and VIX at $T_1$, but the multi-maturity problem introduces an identification question that is absent from any one block.
A common construction solves independent two-maturity problems and stitches them recursively through their shared SPX marginal.
This is computationally attractive, but it completes the otherwise unidentified cross-period dependence by assuming that the next SPX--VIX block is conditionally independent of the earlier path given the shared SPX level.
That assumption cannot be tested with the blockwise calibration instruments themselves.
Expiration misalignment creates a separate practical difficulty because no traded SPX marginal is available at a VIX-only date.
Indeed, \citet[Section~8]{guyon2020joint} observes that a fully coupled treatment across all maturities is ``doable in principle but impractical'' because the problem dimension grows exponentially.

This paper formulates the full-vector treatment and separates two issues that stitching can obscure.
At the exact level, local and global feasibility are equivalent, but the globally admissible path-law class can be strictly larger: a block-preserving Markovian projection leaves every monthly calibration unchanged while discarding cross-period volatility information.
At the numerical level, independently processed surfaces and restricted finite supports can make quote, martingale, and dispersion rows stiff or incompatible at the requested tolerances.
The first issue is one of identification and model risk; the second is one of robust calibration.
Our global formulation provides the common state space needed for both, while the proposed augmented-Bregman mirror-descent scheme addresses the latter by applying cyclic Bregman corrections to observable quote moments and penalizing the conditional residuals.

\subsection{Related work}\label{subsec:related}

Early parametric approaches, double CEV~\citep{gatheral2008consistent}, regime-switching Heston~\citep{papanicolaou2014regime}, and jump diffusions~\citep{cont2013consistent,baldeaux2014consistent}, offer flexibility but generally fail to jointly calibrate at short maturities.
Rough volatility models~\citep{gatheral2018volatility,jacquier2025rough,gatheral2020quadratic} and neural SDEs~\citep{guyon2023neural} improve accuracy at substantial computational cost.
Computationally, the two-maturity MOT problem is often approached with Sinkhorn-type algorithms~\citep{cuturi2013sinkhorn,altschuler2017near,Benamou2015Iterative,lin2022complexity}; conditional martingale and dispersion rows require additional projection or optimization steps, and the state dimension grows quickly with the number of maturities.
The closest exact-calibration benchmark is \citet{bourgey2024fast}, who develop fast discrete- and continuous-time methods for the joint SPX--VIX smile problem.
Alternative routes use signature methods~\citep{cuchiero2025joint} and polynomial diffusions~\citep{abi2025joint}.
Beyond calibration, \citet{che2026spx} compute risk sensitivities within the entropic MOT setting for a \emph{fixed} calibrated coupling; our contribution addresses the antecedent multi-maturity formulation and its finite-dimensional reconciliation.
Our block-preserving construction is also related in spirit to Markovian projection and mimicking~\citep{gyongy1986mimicking,brunick2013mimicking}; the precise distinction is given in \cref{rem:mimicking}.

\subsection{Contributions}\label{subsec:contributions}

\paragraph{Modeling: from stitching to global coupling.}
\begin{enumerate}[label=(\roman*)]
    \item We prove that exact local and global feasibility are equivalent and construct a block-preserving Markovian projection from every global law to a stitched law (\cref{thm:feasibility_equivalence}).
    \item We prove that the stitched full-path law class can nevertheless be a strict subset of the global class: the projection preserves every monthly calibration payoff but can change cross-period dependence and prices (\cref{thm:strict_inclusion,prop:information_loss}).
    \item We show that the standard Markov reference selects the stitched law as the minimum-information exact KL completion (\cref{prop:kl_markovization}). This identifies precisely what must be added (a history-dependent prior, cross-period target, or global payoff objective) to select non-Markov memory.
    \item We characterize compatibility under quote bands: individually feasible relaxed blocks need not admit a common shared-marginal selection, and independent solves can choose different representatives of the shared SPX smile; one global solve enforces a common selection (\cref{prop:band_compatibility,subsec:seam}).
\end{enumerate}

\paragraph{Algorithm: robust finite-dimensional reconciliation.}
\begin{enumerate}[label=(\roman*),resume]
    \item We propose an \emph{augmented-Bregman mirror-descent} scheme that applies cyclic Bregman corrections to normalization, forward, and selected option-price constraints while penalizing the conditional families.
    A finite-grid penalty limit identifies the corresponding least-violating law (\cref{prop:penalty_limit,subsec:augmented}).
    \item On a controlled infeasible affine system, the hard/soft split keeps prescribed marginals about $25$ times tighter than cyclic row projection by leaving a reported conditional residual (\cref{tab:feasibility}).
    \item On smoothed SPX and VIX surfaces, numerical calculations illustrate the practical seam under independent relaxed calibration, the finite-budget quote--conditional tradeoff, and the different numerical roles of the martingale and dispersion penalties (\cref{subsec:market_seam,subsec:market_penalty,subsec:market_projection}).
    \item We extend the formulation to misaligned dates using free SPX nodes and give an exact reference construction for non-overlapping VIX spans (\cref{sec:misaligned,prop:merged_reference}).
\end{enumerate}

\subsection{Organization}\label{subsec:organization}

\Cref{sec:setting} fixes the setting, notation, and the three constraint families.
\Cref{sec:stitching_to_global} develops the modeling contribution: it distinguishes exact stitching from independent relaxed block solves, establishes feasibility equivalence and strict path-law inclusion, and states the no-arbitrage framework.
\Cref{sec:reference} constructs the reference measure and identifies the stitched law selected by the standard entropy criterion.
\Cref{sec:algorithms} develops the finite-dimensional hard/soft method and its penalty-limit interpretation.
\Cref{sec:misaligned} gives the merged-timeline formulation.
\Cref{sec:experiments} collects the finite-state checks and market-data illustrations.
\Cref{sec:conclusion} concludes.

\section{Setting and Notation}\label{sec:setting}

We work in the setting of \citet{guyon2020joint}, extended to $m$ SPX maturities and $m-1$ VIX maturities.
Fix maturities $T_1 < T_2 < \cdots < T_m$ with $T_{i+1} - T_i = \tau = 30/365$ years.
We assume zero interest rates, repos, and dividends for simplicity.
We write $S_i := S_{T_i}$ for the SPX value at $T_i$, $V_i$ for the VIX at $T_i$, and introduce the strictly convex log-contract payoff
\begin{equation}\label{eq:log_contract}
    L(x) := -\frac{2}{\tau}\ln x.
\end{equation}
By definition of the VIX, the price at $T_i$ of the forward-starting log-contract paying $-\frac{2}{\tau}\ln\frac{S_{i+1}}{S_i}$ at $T_{i+1}$ is $V_i^2$.

For probability measures $\mu$ and $\nu$, write
\[
D_{\mathrm{KL}}(\mu\Vert\nu)
:=
\begin{cases}
\displaystyle\int\log\!\left(\frac{d\mu}{d\nu}\right)d\mu,
& \mu\ll\nu,\\[4pt]
+\infty, & \text{otherwise}.
\end{cases}
\]

The market data consists of risk-neutral marginals $\mu_{S_i}$ for $i=1,\ldots,m$ (extracted from SPX option prices at each maturity, e.g.\ via Breeden--Litzenberger~\citep{breeden1978prices}) and $\mu_{V_i}$ for $i=1,\ldots,m-1$ (extracted from VIX futures and option prices).
We denote by $S_0 > 0$ the initial SPX value.

\begin{assumption}\label{ass:marginals}
The given marginals satisfy
\begin{equation}
    \Embb^{\mu_{S_i}}[S_i] = S_0, \quad \Embb^{\mu_{S_i}}[|\ln S_i|] < \infty, \quad i = 1, \ldots, m;
\end{equation}
\begin{equation}
    \Embb^{\mu_{V_i}}[V_i^2] < \infty, \quad i = 1, \ldots, m-1.
\end{equation}
Moreover, the consistency condition
\begin{equation}\label{eq:consistency_global}
    \Embb^{\mu_{V_i}}[V_i^2] = \Embb^{\mu_{S_{i+1}}}[L(S_{i+1})] - \Embb^{\mu_{S_i}}[L(S_i)], \quad i = 1, \ldots, m-1
\end{equation}
holds, and $\mu_{S_i}$, $\mu_{S_{i+1}}$ are in convex order for each $i$.
\end{assumption}

\paragraph{The three constraint families.} Throughout the paper a calibrated model is a law $\mu$ over the merged-timeline state $(S_1, V_1, S_2, V_2, \ldots, V_{m-1}, S_m)$ that satisfies three families of constraints, which we name once and reuse:
\begin{itemize}
    \item \textbf{(C1) Marginal:} each SPX/VIX smile is matched:
    $S_i \sim \mu_{S_i}$ and $V_i \sim \mu_{V_i}$.
    \item \textbf{(C2) Martingale:} the forward-adjusted conditional mean of the next SPX level equals the current one (here $\Embb^\mu[S_{i+1}\mid\cF_i]=S_i$ under zero rates; in general $=\tfrac{F_{i+1}}{F_i}S_i$).
    \item \textbf{(C3) Dispersion:} the conditional log-contract equals $\mathrm{VIX}^2$:
    $\Embb^\mu[L(S_{i+1}/S_i)\mid\cF_i]=V_i^2$.
\end{itemize}
The single structural distinction between the two frameworks studied here is how adjacent blocks are joined.
Within each monthly block, stitching imposes (C2)--(C3) conditional on $(S_i,V_i)$, but it draws the block $(V_i,S_{i+1})$ from the earlier path using only the shared SPX state $S_i$.
The \emph{global} model instead imposes (C2)--(C3) conditional on the entire history $\cF_i=\sigma(S_1,V_1,\ldots,S_i,V_i)$ and does not impose this seam-wise conditional independence.

\begin{remark}[Two notations]\label{rem:notation}
We use the \emph{market-time} notation $(S_1,V_1,S_2,V_2,\ldots,S_m)$ for exposition.
The discretized solver uses an equivalent \emph{solver-coordinate} representation in which each transition is generated by an innovation variable $Z_i$ through $S_{i+1} = f_i(S_i, V_i, Z_i)$; we introduce it only where the discretization is described (\cref{subsec:discrete_formulation}) and otherwise work in market time.
\end{remark}

\section{From Stitching to Global Coupling}\label{sec:stitching_to_global}

This section motivates the global formulation.
After recalling the stitched construction (\cref{subsec:stitching}), we separate a \emph{practical} issue, independently solved relaxed blocks can select different representatives of a shared SPX smile (\cref{subsec:seam}), from a \emph{structural} one, exact stitching imposes a Markov restriction that is invisible to every monthly calibration block (\cref{subsec:stitching_limits}).
We then present the globally coupled model and its duality theory (\cref{subsec:global_model,subsec:duality}) and close by stating the computational difficulty it creates, which the augmented-Bregman solver of \cref{sec:algorithms} resolves (\cref{subsec:handoff}).

\subsection{The Stitching Approach}\label{subsec:stitching}

For each $i = 1, \ldots, m-1$, the monthly feasible set is defined as in \citet{guyon2020joint}.

\begin{definition}[Monthly feasible set]\label{def:monthly_feasible}
$\cP_i := \cP(\mu_{S_i}, \mu_{V_i}, \mu_{S_{i+1}})$ is the set of all probability measures $\nu_i$ on $\mathbb{R}_{>0} \times \mathbb{R}_{\geq 0} \times \mathbb{R}_{>0}$ such that
\begin{gather}\label{eq:monthly_constraints}
    S_i \sim \mu_{S_i}, \quad V_i \sim \mu_{V_i}, \quad S_{i+1} \sim \mu_{S_{i+1}}, \nonumber\\
    \Embb^{\nu_i}[S_{i+1}|S_i, V_i] = S_i, \qquad \Embb^{\nu_i}\!\left[L\!\left(\frac{S_{i+1}}{S_i}\right)\middle| S_i, V_i\right] = V_i^2.
\end{gather}
\end{definition}

\begin{definition}[Full feasible set]\label{def:full_feasible}
Let $\cP_{\mathrm{full}}$ denote the set of probability measures $\mu$ on the state
\[
(S_1,V_1,S_2,V_2,\ldots,V_{m-1},S_m)
\]
satisfying
\begin{align}
    (C1) &\quad S_i \sim \mu_{S_i},\quad i=1,\ldots,m;
    \qquad V_i \sim \mu_{V_i},\quad i=1,\ldots,m-1, \label{eq:C1}\\
    (C2) &\quad \Embb^\mu[S_{i+1}\mid\cF_i]=S_i,
    \quad i=1,\ldots,m-1, \label{eq:C2}\\
    (C3) &\quad \Embb^\mu\!\left[
    L\!\left(\frac{S_{i+1}}{S_i}\right)\middle|\cF_i
    \right]=V_i^2,
    \quad i=1,\ldots,m-1, \label{eq:C3}
\end{align}
where $\cF_i=\sigma(S_1,V_1,\ldots,S_i,V_i)$.
\end{definition}

\begin{definition}[Stitched model]\label{def:stitched}
Given $\nu_i \in \cP_i$ for each $i = 1, \ldots, m-1$, the stitched model $\mu_{\mathrm{stitch}}$ on $\mathbb{R}_{>0}^{m} \times \mathbb{R}_{\geq 0}^{m-1}$ is defined recursively: $(S_1, V_1, S_2) \sim \nu_1$ and $\Law(V_{i+1}, S_{i+2} \mid S_1, V_1, \ldots, S_{i+1}) = \Law_{\nu_{i+1}}(V_{i+1}, S_{i+2} \mid S_{i+1})$ for $i = 1, \ldots, m-2$.
\end{definition}

The critical structural consequence is that the stitched model is \emph{Markovian in the SPX}: the future evolution of $S$ and $V$ depends on the past only through the current SPX value $S_{i+1}$.
All memory of the path, prior VIX levels, prior SPX values, the realized trajectory, is discarded at each stitching boundary.
For each month, \citeauthor{guyon2020joint} builds the calibrating model by solving $\inf_{\nu_i \in \cP_i}D_{\mathrm{KL}}(\nu_i\Vert\bar{\mu}_i)$ with a Sinkhorn-type algorithm; a natural reference uses the lognormal transition kernel
\begin{equation}\label{eq:lognormal_kernel}
    S_{i+1} \mid (S_i = s_i, V_i = v_i) \sim s_i \exp\!\left(v_i\sqrt{\tau}\,G - \tfrac{1}{2}v_i^2\tau\right), \qquad G \sim \mathcal{N}(0,1),
\end{equation}
which automatically satisfies the martingale and dispersion constraints.

\subsection{The Practical Seam under Independent Relaxed Calibration}\label{subsec:seam}

Under \cref{def:monthly_feasible}, adjacent \emph{exact} blocks prescribe the same marginal $\mu_{S_{i+1}}$ and therefore admit a consistent gluing.
There is no marginal-level seam in that idealized definition.
A common relaxed block implementation is different: each two-maturity problem is solved on $(S_i,V_i,Z_i)$ with quote moments imposed on its initial SPX and VIX axes, while the terminal SPX law is the induced transition pushforward and is not separately pinned to the next SPX smile.
A shared maturity $T_{i+1}$ then has two numerical representations, the terminal pushforward of block $i$ and the quote-calibrated initial axis of block $i+1$.
Before those representations are reconciled, the two candidate blocks do not define one exact stitched law.
It is useful to distinguish this practical seam from the structural Markov restriction.

\begin{itemize}
    \item \emph{Practical marginal seam.} The terminal pushforward of block $i$ may differ from the quote-calibrated initial marginal of block $i+1$.
    \item \emph{Structural Markov seam.} Even when the shared marginal is exactly common, stitching draws the next block conditionally only on $S_{i+1}$ and therefore discards dependence on the earlier path.
\end{itemize}

\begin{remark}[One shared representation in the global program]
\label{rem:seam_global}
The globally coupled model of \cref{subsec:global_model} uses a single shared SPX representation at each maturity for both the incoming and the outgoing transition.
It therefore removes the practical marginal seam by construction and does not impose the structural conditional-independence restriction.
\end{remark}

Such independently relaxed blocks need not form one law until their two representations of the shared SPX maturity are reconciled.
A global solve eliminates this duplication by carrying one $S_{i+1}$ variable through both adjacent transitions.
This practical point is separate from exact gluing: when both exact blocks prescribe the same $\mu_{S_{i+1}}$, \cref{thm:feasibility_equivalence} guarantees that they can be glued.

This practical distinction can be stated exactly when quotes are imposed as bands.
Let $\mathfrak M_{S_i}$ and $\mathfrak M_{V_i}$ denote the sets of one-dimensional laws whose selected option moments lie inside the corresponding bid--ask intervals, and define the local feasibility relation
\begin{equation}\label{eq:local_feasibility_relation}
\mathfrak R_i
:=
\left\{(\alpha_i,\beta_i,\alpha_{i+1})
\in\mathfrak M_{S_i}\times\mathfrak M_{V_i}\times\mathfrak M_{S_{i+1}}:
\cP(\alpha_i,\beta_i,\alpha_{i+1})\neq\emptyset\right\}.
\end{equation}
Let $\cP_{\mathrm{full}}^{\mathrm{band}}$ be the set of laws satisfying (C2)--(C3) whose one-dimensional marginals belong to these band-admissible sets.

\begin{proposition}[Compatibility under quote bands]\label{prop:band_compatibility}
The global band-constrained problem, $\cP_{\mathrm{full}}^{\mathrm{band}}\neq\emptyset$, is feasible if and only if there is a \emph{single} sequence of marginal laws
\[
(\alpha_1,\beta_1,\ldots,\beta_{m-1},\alpha_m)
\]
such that $(\alpha_i,\beta_i,\alpha_{i+1})\in\mathfrak R_i$ for every $i$.
Equivalently, the compatible fiber product of the relations $\mathfrak R_i$ over their shared SPX marginals is nonempty.
For $m\geq3$, checking only $\mathfrak R_i\neq\emptyset$ separately is not a certificate of global feasibility, because it does not verify that adjacent relations admit the same shared-marginal selection.
\end{proposition}
The proof is given in \cref{app:proof_band_compatibility}.

Thus a jointly optimized global problem can succeed where a \emph{prescribed} sequential rule fails, for example when that rule fixes an interpolated shared marginal outside the compatibility region even though another marginal inside the quote band is compatible with both neighboring blocks.
It cannot succeed where every admissible version of an exact monthly problem is infeasible.
Likewise, introducing the same conditional tolerances on both sides does not reverse the implication: if $Y_i$ denotes either conditional pricing error, then
\[
\left\|\Embb^\mu[Y_i\mid S_i,V_i]\right\|_{L^p}
\leq
\left\|\Embb^\mu[Y_i\mid\cF_i]\right\|_{L^p},
\qquad 1\leq p\leq\infty,
\]
by conditional Jensen.
The stronger practical claim therefore concerns joint selection and minimum relaxation, not an enlargement of exact feasibility.

\subsection{Structural Limitations of Stitching}\label{subsec:stitching_limits}

The exact distinction is not feasibility but identification of the full path law.
Define
\begin{equation}\label{eq:stitched_class}
\cP_{\mathrm{stitch}}
:=
\left\{
\nu_1(ds_1,dv_1,ds_2)
\prod_{i=2}^{m-1}\nu_i(dv_i,ds_{i+1}\mid s_i)
:\ \nu_i\in\cP_i
\right\}.
\end{equation}
For $\mu\in\cP_{\mathrm{full}}$, let $\nu_i^\mu:=\Law_\mu(S_i,V_i,S_{i+1})$ and define its SPX-Markovization by
\begin{equation}\label{eq:markovization}
\mathsf M\mu
:=
\nu_1^\mu(ds_1,dv_1,ds_2)
\prod_{i=2}^{m-1}\mu(dv_i,ds_{i+1}\mid s_i).
\end{equation}

\begin{theorem}[Feasibility equivalence and block-preserving Markovization]
\label{thm:feasibility_equivalence}
Assume regular conditional distributions exist and adjacent blocks use the same prescribed shared SPX marginal.
For every $\mu\in\cP_{\mathrm{full}}$,
\begin{equation}\label{eq:block_preservation}
    \mathsf M\mu\in\cP_{\mathrm{stitch}}\subseteq\cP_{\mathrm{full}},
    \qquad
    \Law_{\mathsf M\mu}(S_i,V_i,S_{i+1})
    =\Law_\mu(S_i,V_i,S_{i+1})
\end{equation}
for every $i$.
Consequently,
\begin{equation}\label{eq:feasibility_equivalence}
\cP_{\mathrm{full}}\neq\emptyset
\quad\Longleftrightarrow\quad
\cP_{\mathrm{stitch}}\neq\emptyset
\quad\Longleftrightarrow\quad
\cP_i\neq\emptyset\ \text{for every }i.
\end{equation}
\end{theorem}
The proof is given in \cref{app:proof_feasibility_equivalence}.

\begin{remark}[Relation to mimicking]\label{rem:mimicking}
The map $\mu\mapsto\mathsf M\mu$ is a block-valued discrete-time Markovian projection.
It extends the elementary discrete-time Markovization that preserves adjacent pair laws by preserving the overlapping SPX--VIX blocks $(S_i,V_i,S_{i+1})$.
It is conceptually related to continuous-time mimicking results such as those of \citet{gyongy1986mimicking} and \citet{brunick2013mimicking}, but it is not a diffusion-projection theorem: here the result follows directly from disintegration and Markov gluing.
\end{remark}

\begin{theorem}[Strict law-class inclusion and local non-identification]
\label{thm:strict_inclusion}
For $m\geq3$, there exist marginals satisfying \cref{ass:marginals} for which
\begin{equation}
    \cP_{\mathrm{stitch}}\subsetneq\cP_{\mathrm{full}}.
\end{equation}
More precisely, there is a law $\mu\in\cP_{\mathrm{full}}$ with $\mu\neq\mathsf M\mu$ although all identities in~\eqref{eq:block_preservation} hold.
Thus every integrable block-local payoff $f_i(S_i,V_i,S_{i+1})$ has the same price under $\mu$ and $\mathsf M\mu$, while some bounded cross-period payoff $f$ has different prices under the two laws.
Moreover, $\mu=\mathsf M\mu$ if and only if
\begin{equation}\label{eq:stitched_ci}
    (V_i,S_{i+1})\ \perp\ \cF_{i-1}\mid S_i,
    \qquad i=2,\ldots,m-1,
\end{equation}
up to almost-sure equality of the corresponding kernels.
\end{theorem}
The proof and finite-state construction are given in \cref{app:proof_strict_inclusion}.

\begin{proposition}[Information removed by stitching]\label{prop:information_loss}
Let $\mu\in\cP_{\mathrm{full}}$ and write the pre-$S_i$ history as
\[
H_{i-1}:=(S_1,V_1,\ldots,S_{i-1},V_{i-1}).
\]
Suppose $D_{\mathrm{KL}}(\mu\Vert\mathsf M\mu)<\infty$.
Then
\begin{equation}\label{eq:information_loss}
D_{\mathrm{KL}}(\mu\Vert\mathsf M\mu)
=
\sum_{i=2}^{m-1}
I_\mu\!\left((V_i,S_{i+1});H_{i-1}\mid S_i\right).
\end{equation}
In particular, the information loss is zero exactly when the stitched conditional-independence relations~\eqref{eq:stitched_ci} hold.
\end{proposition}
The proof is given in \cref{app:proof_information_loss}.

The implication for calibration is a negative-identification result.
No collection of payoffs measurable within the individual blocks can test the Markov restriction, because $\mu$ and $\mathsf M\mu$ agree on every such payoff.
Stitching therefore does more than assemble local solutions: it completes the unidentified temporal dependence by imposing~\eqref{eq:stitched_ci}.
This can alter VIX persistence, joint tails, stress exposures, hedge sensitivities, and prices of products spanning several volatility windows.

The exact finite-tree construction illustrates the effect.
Markovization attenuates $\Cov(V_1,V_2)$ from $0.040$ to $0.011$, a $72\%$ reduction, and changes the price of the VIX-spread call $(V_2-V_1-K)^+$ from $(0.050,0,0,0)$ to $(0.113,0.063,0.045,0.027)$ at strikes $K=(0,0.1,0.2,0.3)$.
Both laws satisfy the same local constraints and have the same adjacent triple laws; the price difference is entirely due to the cross-period gluing.

More generally, for compatible local laws $\nu_1,\ldots,\nu_{m-1}$ define the set of global gluings
\begin{equation}\label{eq:gluing_class}
\Gamma(\nu_1,\ldots,\nu_{m-1})
:=
\left\{\mu\in\cP_{\mathrm{full}}:
\Law_\mu(S_i,V_i,S_{i+1})=\nu_i,
\ i=1,\ldots,m-1\right\}.
\end{equation}
For a cross-period payoff $f$, the interval
\begin{equation}\label{eq:robust_bounds}
\underline p(f):=\inf_{\mu\in\Gamma}\Embb^\mu[f],
\qquad
\overline p(f):=\sup_{\mu\in\Gamma}\Embb^\mu[f]
\end{equation}
measures dependence risk left unidentified by all monthly calibration instruments.
The stitched law supplies one point in this interval, not a model-independent price.
Computing such bounds for market portfolios is left to future work; the finite-tree example establishes that the interval can be non-degenerate.

\begin{figure}[ht]
\centering
\includegraphics[width=\textwidth]{figures/p1_feasibility.pdf}
\caption{Block-preserving Markovization on the exact finite-tree example of \cref{thm:strict_inclusion}.
The global law and its stitched Markovization have identical adjacent calibration blocks, but the covariance $\Cov(V_1,V_2)$ falls from $0.040$ to $0.011$ and the prices of VIX-spread calls change.
The discrepancy is a pure loss of cross-period information, not a feasibility difference.}
\label{fig:markov_gap}
\end{figure}

\subsection{The Globally Coupled Model}\label{subsec:global_model}

The global model calibrates one law over the entire chain and imposes the martingale and dispersion identities with respect to the full history $\cF_i$ in \cref{def:full_feasible}.
It does not impose the seam-wise conditional independence~\eqref{eq:stitched_ci}; the stitched class is recovered as the subset generated by the Markovization~\eqref{eq:markovization}.
This additional path-law freedom does not change exact local feasibility, but it is necessary for representing cross-period information and for computing dependence-sensitive model-risk bounds.
Let $\Omega$ denote the admissible path space, or the allowed set of cells in a finite discretization.

\begin{definition}[Full trading universe]\label{def:full_trading}
Let $\cU_{\mathrm{full}}$ denote the set of all portfolios
\[
\Pi = \bigl(\{u_{S_i}\}_{i=1}^{m},\; \{u_{V_i}\}_{i=1}^{m-1},\; \{\Delta_{S,i}\}_{i=1}^{m-1},\; \{\Delta_{L,i}\}_{i=1}^{m-1}\bigr)
\]
where $u_{S_i}$, $u_{V_i}$ are integrable vanilla payoffs and $\Delta_{S,i}$, $\Delta_{L,i}$ are bounded $\cF_i$-measurable dynamic trading strategies.
Its initial cost is
\[
\mathrm{Cost}(\Pi)
:=
\sum_{i=1}^{m}\int u_{S_i}\,d\mu_{S_i}
+\sum_{i=1}^{m-1}\int u_{V_i}\,d\mu_{V_i},
\]
because the dynamic gains have zero initial cost.
The portfolio terminal value is
\begin{multline}\label{eq:portfolio_value}
    \Pi(\bs, \bv) = \sum_{i=1}^{m} u_{S_i}(s_i) + \sum_{i=1}^{m-1} u_{V_i}(v_i) \\
    + \sum_{i=1}^{m-1} \Delta_{S,i}(\bs^i, \bv^i)(s_{i+1} - s_i) + \sum_{i=1}^{m-1} \Delta_{L,i}(\bs^i, \bv^i)\!\left(L\!\left(\frac{s_{i+1}}{s_i}\right) - v_i^2\right).
\end{multline}
A \emph{global $(S,V)$-arbitrage} is a portfolio $\Pi \in \cU_{\mathrm{full}}$ with $\Pi(\bs,\bv) \geq 0$ for every $(\bs,\bv)\in\Omega$ and $\mathrm{Cost}(\Pi) < 0$.
\end{definition}

Given a reference measure $\bar{\mu}$, the entropic calibration problem is
\begin{equation}\label{eq:entropic_full}
    D_{\bar{\mu}}^{\mathrm{full}} := \inf_{\mu \in \cP_{\mathrm{full}}}D_{\mathrm{KL}}(\mu\Vert\bar{\mu}).
\end{equation}
On a finite grid and under the usual constraint qualification, the optimizer has the normalized exponential form
\begin{multline}
\frac{d\mu^*}{d\bar\mu}(\bs,\bv)
=\frac{1}{Z}\exp\!\Biggl(
\sum_i u^*_{S_i}(s_i)+\sum_i u^*_{V_i}(v_i)\\
+\sum_i\Delta^*_{S,i}(\bs^i,\bv^i)(s_{i+1}-s_i)
+\sum_i\Delta^*_{L,i}(\bs^i,\bv^i)
\left[L(s_{i+1}/s_i)-v_i^2\right]
\Biggr),
\end{multline}
where $Z$ normalizes the law and the dynamic multipliers may depend on the full history.

\subsection{Finite-State Duality and Arbitrage-Freeness}\label{subsec:duality}

The easy direction of the no-arbitrage statement is immediate: any $\mu\in\cP_{\mathrm{full}}$ prices every portfolio in \cref{def:full_trading} at its initial cost and therefore rules out the stated arbitrage.
We state the converse on the finite state spaces used by the numerical checks, where it is a direct alternative theorem for linear systems.
Extending that converse to an unbounded continuum path space requires a separate specification of admissible payoff growth, topology, and closedness; we do not claim that extension here.

\begin{theorem}[Finite-state arbitrage duality]\label{thm:duality}
Every law in $\cP_{\mathrm{full}}$ rules out global $(S,V)$-arbitrage.
Conversely, on a finite state space, if the trading universe contains the static and dynamic payoffs dual to every row defining (C1)--(C3), absence of global $(S,V)$-arbitrage implies $\cP_{\mathrm{full}}\neq\emptyset$.
\end{theorem}
The proof is given in \cref{app:proof_duality}.

\begin{theorem}[The entropic model is globally arbitrage-free]\label{thm:entropic_af}
Suppose $\cP_{\mathrm{full}}$ is weakly closed and there exists $\mu_0\in\cP_{\mathrm{full}}$ with $D_{\mathrm{KL}}(\mu_0\Vert\bar\mu)<\infty$.
If a relative-entropy sublevel containing a minimizing sequence is weakly compact, then the minimizer
$\mu^*=\argmin_{\mu\in\cP_{\mathrm{full}}}D_{\mathrm{KL}}(\mu\Vert\bar\mu)$ exists and is unique.
As a feasible law, it rules out global arbitrage.
On the finite grids used below, these existence conditions reduce to positivity of the reference on the allowed cells and nonemptiness of the exact constraint set.
\end{theorem}
The proof is given in \cref{app:proof_entropic_af}.

The chain of implications is
\begin{equation}\label{eq:chain}
    \cP_{\mathrm{stitch}} \neq \emptyset
    \;\;\Longleftrightarrow\;\;
    \cP_{\mathrm{full}} \neq \emptyset
    \;\;\Longleftrightarrow\;\;
    \text{No global arbitrage}
\end{equation}
on a finite state space under the payoff-completeness condition of \cref{thm:duality}.
The equality concerns nonemptiness, not equality of law classes; \cref{thm:strict_inclusion} shows that the latter can be strict.

\begin{table}[ht]
\centering
\caption{Structural comparison of the stitching and globally coupled approaches.}
\label{tab:comparison}
\small
\begin{tabular}{@{}p{4.2cm}p{4.8cm}p{4.8cm}@{}}
\toprule
\textbf{Aspect} & \textbf{Stitching} & \textbf{Globally Coupled} \\
\midrule
Dimensionality per subproblem & 3-dimensional $(S_i, V_i, S_{i+1})$ & $(2m{-}1)$-dimensional full vector \\
\addlinespace
Markov structure & Conditional independence at each SPX seam & No seam-wise Markov assumption \\
\addlinespace
Delta strategies at $T_i$ & $\Delta(s_i, v_i)$ & $\Delta(\cF_i)$ (full history) \\
\addlinespace
Shared SPX maturity & Common in the exact gluing; duplicated in independent relaxed solves & Single shared representation \\
\addlinespace
Joint $(V_i, V_j)$ distribution & Indirect: $V_i \to S_{i+1} \to V_{i+1}$ & Directly modeled \\
\addlinespace
Exact nonemptiness & Equivalent to global nonemptiness & Equivalent to all monthly problems being nonempty \\
\addlinespace
Maturity misalignment & Requires interpolation or endogenous boundary laws & Supports free nodes; overlapping windows require additional care \\
\bottomrule
\end{tabular}
\end{table}

\subsection{From Modeling to Algorithm: a Single Coupling of Competing Constraints}\label{subsec:handoff}

The global formulation can expose a strictly larger class of path laws (\cref{thm:strict_inclusion,prop:information_loss}) and carries one shared representation through the full chain (\cref{subsec:seam,rem:seam_global}).
It is the appropriate state space for cross-period targets, dependence-sensitive price bounds, and free SPX nodes at misaligned dates.
Its computational price is that all quote, martingale, and dispersion rows act on one high-dimensional tensor.

On finite supports, discretization, quadrature restrictions, and active-row choices can create substantial constraint tension even though exact continuum nonemptiness is equivalent.
The finite-state checks do not claim a general impossibility of exact finite-grid calibration; instead they isolate how the tested projection schemes allocate residuals when a finite affine system is incompatible.
The hard/soft method of \cref{sec:algorithms} assigns priority to observable quote moments and reports the remaining structural residual explicitly.

\section{The Globally Coupled Reference Measure}\label{sec:reference}

The entropic formulation~\eqref{eq:entropic_full} requires a reference measure $\bar{\mu}$ and a feasible law of finite relative entropy with respect to it.
A product-of-marginals reference satisfies the input marginals but generally violates the conditional constraints.
Conversely, a Markov transition reference can satisfy the conditional identities at initialization but does not encode non-Markov temporal information.
We build $\bar\mu$ so that (C2)--(C3) hold at initialization; later quote-row corrections can disturb them, so the numerical method continues to reconcile both families.

\paragraph{Construction.} The reference interleaves, chronologically, lognormal SPX transition kernels with independent draws of the VIX marginals:
\begin{equation}\label{eq:complete_ref}
    \bar{\mu}(d\mathbf{s}, d\mathbf{v}) = \mu_{S_1}(ds_1) \prod_{i=1}^{m-1}\!\Big[\mu_{V_i}(dv_i) \cdot T_i(s_i, v_i, ds_{i+1})\Big],
\end{equation}
where $T_i$ is the lognormal kernel~\eqref{eq:lognormal_kernel}.
The kernel makes $\Embb^{\bar\mu}[S_{i+1}\mid\cF_i]=S_i$ and $\Embb^{\bar\mu}[L(S_{i+1}/S_i)\mid\cF_i]=V_i^2$ hold exactly at initialization, and the independent draws give $V_i\sim\mu_{V_i}$ under $\bar\mu$.
The subsequent marginal corrections generally alter more than one axis, so all calibrated marginal rows must remain in the solver.

\begin{proposition}[Reference-measure properties]\label{prop:ref_props}
The measure $\bar\mu$ in~\eqref{eq:complete_ref} is well-defined whenever the displayed factors are probability measures, satisfies (C2)--(C3), and matches the $S_1$ and VIX marginals at initialization.
Its support is the support induced by $\mu_{S_1}$, the VIX marginals, and the transition kernels.
On a finite grid, positivity of every allowed reference cell implies that every law on that grid is dominated by $\bar\mu$; in the continuum, existence of a feasible finite-entropy law must be assumed separately.
\end{proposition}
The proof is given in \cref{app:proof_entropic_af}.

\begin{proposition}[The standard reference selects the stitched completion]
\label{prop:kl_markovization}
Suppose $\bar\mu$ factorizes through the shared SPX states as in~\eqref{eq:complete_ref}, and let $\mu\in\cP_{\mathrm{full}}$ have finite relative entropy with respect to $\bar\mu$.
Then
\begin{equation}\label{eq:kl_markovization}
D_{\mathrm{KL}}(\mu\Vert\bar\mu)
=
D_{\mathrm{KL}}(\mathsf M\mu\Vert\bar\mu)
+D_{\mathrm{KL}}(\mu\Vert\mathsf M\mu).
\end{equation}
Consequently, every finite-entropy minimizer over $\cP_{\mathrm{full}}$ belongs to $\cP_{\mathrm{stitch}}$.
\end{proposition}
The proof is given in \cref{app:proof_kl_markovization}.

Thus the current reference gives stitching a precise interpretation: it is the minimum-information exact completion of the local calibrations.
The global state space becomes essential when one adds a history-dependent reference, a cross-period target, or a dependence-sensitive payoff objective; it also exposes the range of path laws that the local smiles leave unidentified.

For numerical stability, the reference tensor is assembled in the log domain by iterated outer sums of the log-marginals and log-kernels.

\section{Robust Finite-Dimensional Calibration}\label{sec:algorithms}

We now solve the globally coupled program.
\Cref{subsec:discrete_formulation} states the finite-dimensional problem and the hard/soft split.
\Cref{subsec:priority_test} gives a controlled feasible/infeasible test of that split, and \cref{subsec:augmented} introduces the augmented-Bregman solver.
\Cref{subsec:unification} relates the idealized exact-projection limit to the implemented finite-budget method.

\subsection{Finite-Dimensional Program and the Hard/Soft Split}\label{subsec:discrete_formulation}

Fix $m$ maturities, an initial SPX grid with $n_S$ points, VIX grids with $n_V$ points, and innovation grids with $n_Z$ points.
We discretize the law as a tensor $\pi$ on innovation coordinates
\[
(S_1,V_1,Z_1,V_2,Z_2,\ldots,V_{m-1},Z_{m-1}),
\]
with
$S_{i+1}=f_i(S_i,V_i,Z_i)$ derived recursively.
For common axis sizes its dimension is
$d=n_S n_V^{m-1}n_Z^{m-1}$; the corresponding law in market coordinates is its pushforward.

The discrete problem matches finitely many forward and call-price moments rather than entire continuous marginals.
Writing all selected rows compactly as $A\pi=b$ with $\pi$ on the probability simplex $\Delta_d$, the ideal exact discrete problem is
\begin{equation}\label{eq:Pm}
    \min_{\pi\in\Delta_d}\; \sum_{z} \pi(z)\,\ln\frac{\pi(z)}{\bar\pi(z)}
    \quad\text{s.t.}\quad A\pi = b,
\end{equation}
where $A$ stacks the quote-moment rows approximating (C1), the conditional martingale rows (C2), and the conditional dispersion rows (C3).
We never materialize $A$; all operations act directly on $\pi$.

\paragraph{Conditional-residual metrics.} For a finite-grid implementation, conditional accuracy can be summarized by relative, mass-weighted residuals over populated conditioning cells.
Fix a transition $i$ and a conditioning cell $c$ (an atom of the filtration $\cF_i$); write $w_c$ for its probability mass and $\Embb[\,\cdot\mid\cF_i{=}c\,]$ for the corresponding conditional expectation.
The exact theory above is written in forward units.
When the coordinates retain nonconstant market forwards, let $d_i:=F_{i+1}/F_i$ and evaluate both conditions on the forward-adjusted transition.
The per-cell relative residuals are
\begin{equation}\label{eq:rel_residuals}
    r^{\mathrm{disp}}_{i,c} = \frac{\bigl|\,\Embb[L(S_{i+1}/(d_iS_i))\mid\cF_i{=}c] - V_i^2\,\bigr|}{V_i^2},
    \qquad
    r^{\mathrm{mart}}_{i,c} = \left|\,\frac{\Embb[S_{i+1}\mid\cF_i{=}c]}{d_iS_i} - 1\,\right|,
\end{equation}
the dispersion residual measured relative to $\mathrm{VIX}_i^2$ and the martingale residual relative to the forward-adjusted level.
For each family we report the mass-weighted mean of these residuals over the cells carrying at least $10\%$ of the peak conditioning mass, worst-cased over transitions:
\begin{equation}\label{eq:massw}
    \edisp \;=\; \max_i \!\!\sum_{c:\,w_c\ge 0.1\,w_{\max}}\!\! \tilde w_c\, r^{\mathrm{disp}}_{i,c},
    \qquad
    \emart \;=\; \max_i \!\!\sum_{c:\,w_c\ge 0.1\,w_{\max}}\!\! \tilde w_c\, r^{\mathrm{mart}}_{i,c},
\end{equation}
with $\tilde w_c$ the masses renormalized over the retained cells.
The displayed $10\%$ threshold is illustrative: any thresholded statistic must be accompanied by its retained probability mass and an all-cell statistic before it is used to discuss feasibility.

\paragraph{Constraint treatment.} The constraint families do not play symmetric computational roles, and our solver treats them differently:
\begin{itemize}
    \item \textbf{Bregman-corrected:} normalization, forward, and selected call-price constraints are visited cyclically with clipped Newton/Bregman corrections at every inner sweep.
    \item \textbf{Penalized:} the active conditional martingale and dispersion rows enter the mirror gradient with tunable weights.
\end{itemize}
Each row correction targets its affine hyperplane, but one finite cyclic pass need not leave the iterate in the intersection of all quote rows.
Accordingly, we report the final quote and implied-volatility errors rather than calling the finite-budget iterates exactly marginal-feasible.
The rationale for the split is economic: observable option prices receive priority, while discrepancies caused by input or finite-support tension remain visible in the structural residuals.

To formalize this treatment independently of a particular iteration schedule, let
$A_h\pi=b_h$ collect the normalization, forward, and selected option-price equalities, let
$\mathcal H:=\{\pi\in\Delta_d:A_h\pi=b_h\}$ be nonempty, and let $A_c\pi=b_c$ collect the penalized conditional rows.
Set
\begin{equation}\label{eq:penalty_residual}
r(\pi):=\tfrac12\lVert A_c\pi-b_c\rVert^2,
\end{equation}
which gives one unscaled conditional residual model.
For $\lambda>0$, consider
\begin{equation}\label{eq:penalty_problem}
\pi_\lambda
:=
\argmin_{\pi\in\mathcal H}
\left\{D_{\mathrm{KL}}(\pi\Vert\bar\pi)+\lambda r(\pi)\right\}.
\end{equation}

\begin{proposition}[Penalty limit on a finite grid]\label{prop:penalty_limit}
Suppose $\bar\pi$ is positive on the allowed finite support and $\mathcal H$ is nonempty.
Then~\eqref{eq:penalty_problem} has a unique solution.
As $\lambda\to\infty$, $r(\pi_\lambda)$ converges to
$r_\star:=\min_{\pi\in\mathcal H}r(\pi)$, and $\pi_\lambda$ converges to the minimum-KL element of $\argmin_{\pi\in\mathcal H}r(\pi)$.
If the full discrete system is feasible, then $r_\star=0$ and the limit is the exact constrained KL projection; otherwise the limit is the least-violating law that preserves the $A_h$ constraints.
\end{proposition}
The proof is given in \cref{app:proof_penalty_limit}.

The iterative method below uses the same residual gradient inside an augmented-Lagrangian mirror iteration, rather than solving~\eqref{eq:penalty_problem} to convergence for every $\lambda$.
Thus \cref{prop:penalty_limit} supplies an idealized variational interpretation, not a convergence claim for a fixed iteration budget.

\subsection{Controlled Priority-Allocation Test}\label{subsec:priority_test}

To isolate the effect of the hard/soft split from market-data conventions, we use a finite Gaussian martingale-coupling problem with one analytic feasibility knob.
Writing $\rho_{\mathrm{var}}:=\operatorname{Var}(Y)/\operatorname{Var}(X)$, the source and target marginals have the same mean and satisfy $\rho_{\mathrm{var}}\approx1.3$ in the feasible case and $\rho_{\mathrm{var}}\approx0.7$ in the infeasible case, as confirmed by the direct LP check below.
The three methods use the same $41\times41$ state grid and constraints but assign the marginal and conditional rows differently.
The complete specification and results are given in \cref{sec:experiments,app:omd_theory,tab:feasibility}.
On the infeasible instance, cyclic row projection leaves a marginal residual of $1.9\times10^{-2}$, whereas the marginal-priority scheme leaves $7.4\times10^{-4}$, about $25$ times smaller, while exposing a conditional residual of $1.1\times10^{-3}$.
This controlled result supports assigning priority to prescribed marginal rows; it does not establish infeasibility of any empirical SPX--VIX dataset.

\subsection{The Augmented-Bregman Solver}\label{subsec:augmented}

We treat the conditional rows through an augmented-Lagrangian mirror step and interleave it with cyclic Bregman corrections of the $A_h$ rows.
Recall that \emph{mirror descent} is the standard first-order scheme that, at each step, moves along the (sub)gradient of the objective in the geometry induced by a strictly convex mirror map and then projects back with the associated Bregman divergence~\citep{nemirovski1983problem,beck2003mirror}; with the negative-entropy mirror map on the simplex its update is the multiplicative (exponentiated-gradient) reweighting $\pi\odot\exp(-\alpha g)$ followed by renormalization, and the exact KL/Bregman projection onto a linear constraint reduces to the Sinkhorn-type rescaling used for the hard families~\citep{Benamou2015Iterative}.
Let $A_c,b_c$ collect the active conditional rows, with dual variables $y$ and penalty weight $\lambda$.
At inner iterate $\pi_t$, the method forms
\begin{equation}\label{eq:aug_gradient}
g_t=A_c^\top\!\left[y+\lambda(A_c\pi_t-b_c)\right]
\end{equation}
and takes the entropic mirror step
\begin{equation}\label{eq:mirror_step}
\widetilde\pi_{t+1}
=\argmin_{\pi\in\Delta_d}
\left\{\langle g_t,\pi\rangle+\eta_t^{-1}D_{\mathrm{KL}}(\pi\Vert\pi_t)\right\}
\propto \pi_t\odot\exp(-\eta_t g_t),
\end{equation}
where $\eta_t$ is clipped to bound the largest log-weight update.
It then applies one clipped Newton correction in KL geometry to each $A_h$ row.
After the inner budget, the method updates
$y\leftarrow y+\lambda(A_c\pi-b_c)$ and optionally increases $\lambda$.
Because both blocks are revisited throughout the run, the method avoids making one large conditional projection whose entire correction is handed back at the next quote pass.

\begin{algorithm}[H]
\caption{Augmented-Bregman solver for SPX--VIX calibration}
\label{alg:augmented}
\DontPrintSemicolon
\SetKwInOut{Input}{Input}\SetKwInOut{Output}{Output}
\Input{Reference and rows; penalty schedule; $\alpha$, $c$; iteration counts}
\Output{Calibrated tensor $\pi$}
\BlankLine
$\pi \leftarrow \bar\mu$;\quad $y \leftarrow 0$;\quad $\lambda \leftarrow \lambda_0$\;
\For{$k = 1, \ldots, K_{\mathrm{out}}$}{
  \For{$t = 1, \ldots, K_{\mathrm{in}}$}{
    $g \leftarrow A_c^\top\big(y + \lambda (A_c\pi - b_c)\big)$ \tcp*{penalty gradient (sparse)}
        $\eta \leftarrow \min\{\alpha,c/\lVert g\rVert_\infty\}$;\quad $\pi \leftarrow \pi \odot \exp(-\eta g)$;\quad renormalize\;
    Apply one Bregman--Newton correction to each row of $(A_h,b_h)$\;
  }
  $y \leftarrow y + \lambda (A_c \pi - b_c)$;\quad $\lambda \leftarrow \min(\lambda\gamma, \lambda_{\max})$\;
}
\Return $\pi$\;
\end{algorithm}

\subsection{Relation to Exact Projection}\label{subsec:unification}

At the level of exactly solved optimization problems, the quadratic penalty
$\tfrac{\lambda}{2}\lVert A_c\pi-b_c\rVert^2$ converges as $\lambda\to\infty$ to the indicator of the conditional constraint set, provided that set intersects $\{\pi:A_h\pi=b_h\}$.
This idealized limit does not imply that a fixed-budget implementation converges as $\lambda$ grows.
Separately, replacing each first-order mirror step by the exact KL projection onto the visited row recovers cyclic Bregman/Sinkhorn updates.
Thus the methods share an entropic geometry, but finite-budget iterates should not be identified with exact points on an asymptotic solution path (\cref{prop:sinkhorn_omd}).
The mirror-descent view and controlled infeasibility analysis are collected in \cref{app:omd_theory}; the per-sweep cost accounting is in \cref{app:complexity}.

\section{Extension to Misaligned SPX and VIX Maturities}\label{sec:misaligned}

In practice, the SPX and VIX maturities selected for a calibration need not coincide.
When the calibration instruments do not supply an SPX marginal at a VIX settlement date, a standard fixed-marginal monthly block must either interpolate that missing law or introduce it as an endogenous variable.
The global formulation provides a natural framework on a \emph{merged timeline} $\mathcal{T} = \mathcal{T}_S \cup \mathcal{T}_V = \{t_1 < \cdots < t_N\}$, where $\mathcal{T}_S$ are SPX expiries and $\mathcal{T}_V$ are VIX settlement dates, treating the SPX level at VIX-only dates $\mathcal{T}_{\mathrm{free}} := \mathcal{T}_V \setminus \mathcal{T}_S$ as a \emph{free} variable (no marginal constraint there).
The state is $\mathbf{X} = (S_{t_1}, \ldots, S_{t_N}, V_1, \ldots, V_n)$ with filtration $\cF_k := \sigma(\{S_{t_\ell} : t_\ell \leq t_k\} \cup \{V_i :
T_i^V \leq t_k\})$.

\begin{assumption}[Non-overlapping VIX spans]\label{ass:non_overlap}
For each $i$, let $T_{\pi(i)}^S$ be the SPX maturity referenced by the $i$-th VIX contract and $\tau_i := T_{\pi(i)}^S - T_i^V$.
We assume $T_{i+1}^V \geq T_{\pi(i)}^S$ (non-overlapping spans) and $t_1 \in \mathcal{T}_S$.
\end{assumption}

Under \cref{ass:non_overlap}, the feasible set and reference measure extend with three changes.
There is no marginal constraint at $\mathcal{T}_{\mathrm{free}}$; each dispersion constraint spans the merged-grid steps in $[T_i^V,T_{\pi(i)}^S]$; and the filtration at a VIX date includes $V_i$.
The following proposition records the reference construction used in the non-overlapping case.

\begin{proposition}[Non-overlapping merged-timeline reference]\label{prop:merged_reference}
Draw $S_{t_1}$ from a prescribed initial law and draw each $V_i$ from its prescribed marginal at time $T_i^V$, independently of the preceding reference history.
On every merged-grid subinterval contained in $[T_i^V,T_{\pi(i)}^S]$, use a conditionally independent lognormal SPX increment with volatility $V_i$ and the subinterval length; on uncovered subintervals, use any strictly positive martingale transition kernel.
If the VIX spans do not overlap, the chronological product of these kernels is a well-defined martingale reference that matches the initial SPX and all VIX marginals and satisfies
\[
\Embb\!\left[-\frac{2}{\tau_i}\log\!\left(
\frac{S_{T_{\pi(i)}^S}}{d_iS_{T_i^V}}
\right)\middle|\cF_{T_i^V}\right]=V_i^2,
\]
where $d_i$ is the forward ratio over the span.
\end{proposition}
The proof is given in \cref{app:proof_merged_reference}.

The solver of \cref{sec:algorithms} then applies with multi-step dispersion rows.
The aligned case is recovered when $\mathcal{T}_{\mathrm{free}} = \emptyset$.
The definitions and tower-property arguments extend to this timeline; a continuum no-arbitrage converse would require the additional functional-analytic assumptions noted after \cref{thm:duality}.

\begin{proposition}[Standard fixed-marginal stitching is not identified at VIX-only dates]\label{prop:stitch_illposed}
If $T_i^V \notin \mathcal{T}_S$, the supplied SPX vanilla marginals do not determine the law $\mu_{S_{T_i^V}}$ required by a standard three-variable block beginning at $T_i^V$.
Consequently, such a block requires either an exogenous interpolation rule or an enlarged formulation that treats the missing SPX law as endogenous.
\end{proposition}
The proof is given in \cref{app:proof_stitch_illposed}.

The full merged-timeline feasible set and its duality mirror \cref{def:full_feasible,thm:duality} with the three changes above; we omit the restatement.

\paragraph{Overlapping windows.}
If two VIX spans overlap, assigning one volatility state to their shared subinterval generally prevents the product reference from satisfying both dispersion identities automatically.
A consistent extension must model the variance allocated to the overlap and impose martingale rows at every adjacent event time.
We leave the corresponding exact reference construction for future work.
\Cref{subsec:market_misaligned} reports an overlapping-window calibration as a finite-grid numerical comparison; it is not covered by \cref{prop:merged_reference} and is not used to support any feasibility theorem.

\section{Finite-State Checks and Market-Data Illustrations}\label{sec:experiments}

We first report two self-contained finite-state checks whose role is to verify the structural identities and the allocation of residuals across constraint families.
We then report numerical calculations on smoothed SPX and VIX surfaces.
Those market-data calculations illustrate finite-budget calibration behavior; they do not establish a separation between exact local and global feasibility, which is ruled out by \cref{thm:feasibility_equivalence}.

\paragraph{Block-preserving Markovization.}
The finite tree used in \cref{thm:strict_inclusion} is evaluated directly, without numerical optimization.
Its global feasibility residual is $1.42\times10^{-14}$.
Markovization preserves both adjacent triple laws but changes $\Cov(V_1,V_2)$ from $0.0400$ to $0.0110$, a $72.4\%$ attenuation.
For strikes $K=(0,0.1,0.2,0.3)$, the prices of $(V_2-V_1-K)^+$ change from $(0.050,0,0,0)$ to $(0.113,0.063,0.045,0.027)$.

\paragraph{Controlled feasibility knob.}
Let $x_j=0.25+0.0375j$, $j=0,\ldots,40$, and let the source weights be proportional to $\exp(-(x_j-1)^2/(2\sigma_X^2))$ with $\sigma_X=0.15$.
The target weights use $\sigma_Y=\sigma_X\sqrt{1\pm0.30}$ on the same grid; after normalization the discrete variance ratios are $\rho_{\mathrm{var}}\approx1.3$ and $\rho_{\mathrm{var}}\approx0.7$.
We impose both marginals and the martingale rows $\sum_k\pi_{jk}(x_k-x_j)=0$ for the $25$ source states whose marginal mass is at least $1\%$ of the peak.
The cyclic-projection and marginal-priority methods use $800$ sweeps, while all-soft OMD uses $6000$; the soft updates use penalty $200$, base step $0.5$, and log-tilt cap $1$.
The resulting affine systems have $1681$ variables and $107$ rows.
A direct linear-programming check finds the first feasible with maximum equality residual $1.5\times10^{-16}$.
For the contracted case, let $A=\{8,\ldots,32\}$ be the active source indices.
Any feasible coupling would satisfy, by conditional Jensen on the active rows,
\begin{equation}\label{eq:synthetic_infeasibility_certificate}
\sum_k\nu_k(x_k-1)^2
\geq
\sum_{j\in A}\mu_j\bigl(\Embb[Y-1\mid X=x_j]\bigr)^2
=\sum_{j\in A}\mu_j(x_j-1)^2.
\end{equation}
Numerically, the right side is $0.022043550052$, whereas the target variance on the left is $0.015749999517$.
Thus the thresholded affine system is infeasible even though the tail martingale rows are omitted; an independent LP returns the same conclusion.
The entries in \cref{tab:feasibility} are the median maximum residuals over the final $10\%$ of each run and show explicitly how each method allocates the incompatibility.

\subsection{Market-Data Setup and Scope}\label{subsec:market_setup}

We report calculations on P-spline-smoothed SPX and VIX surfaces from the 2026-04-21 snapshot.
These are the P-spline input surfaces used throughout the 2026-04-21 experiment suite.
The five surfaces comprise SPX maturities at $23$, $57$, and $87$ days and VIX maturities at $27$ and $56$ days.
The five-axis implementation pairs the two VIX surfaces with consecutive SPX transition blocks on the solver grid
$(S_1,V_1,Z_1,V_2,Z_2)$ and uses a 30-day log-contract horizon for each block.
Because the actual SPX and VIX event dates do not coincide, this is an aligned-coordinate approximation rather than an exact event-time model of the overlapping windows.

The numerical grid has shape
$(30,25,8,25,8)$, or $1.2\times10^6$ cells.
Selected forward and call-price constraints are revisited through cyclic Bregman corrections, while active martingale and dispersion rows are penalized.
The conditional-row builder activates conditioning cells whose prior mass is at least $1\%$ of the peak prior conditioning mass.
The headline conditional statistics $\edisp$ and $\emart$ are the thresholded bulk metrics of~\eqref{eq:massw}, evaluated on cells whose final conditioning mass is at least $10\%$ of the peak and then worst-cased over the two transitions.
They are finite-grid diagnostics, not certificates that every conditional row is satisfied.
Every penalty value below is an independent fixed-budget run, so the path describes tested operating points rather than exact solutions of~\eqref{eq:penalty_problem}.

These calculations address the numerical question developed in \cref{subsec:handoff,sec:algorithms}: how a finite implementation allocates discrepancies among quote and conditional rows.
They do not show that global coupling rescues an infeasible exact monthly block, nor are they used to prove \cref{thm:feasibility_equivalence,thm:strict_inclusion}.

\subsection{Independent Relaxed Blocks at a Shared Maturity}\label{subsec:market_seam}

The first calculation illustrates the practical seam of \cref{subsec:seam}.
Two relaxed three-variable blocks were calibrated independently and their two numerical representations of the shared $57$-day SPX smile were compared.
At the matched operating point displayed in \cref{fig:market_seam}, the maximum and mean differences over the invertible strike range were approximately $1.1$ and $0.3$ volatility points, respectively.
A separate post-projection ladder, reported in \cref{tab:seam_tightness}, tightened the conditional families in both blocks and found that the shared-smile discrepancy could increase substantially.
The post-projections also change the quote fit, so this calculation concerns independently processed relaxed blocks; it is not a counterexample to the exact gluing result.

\begin{figure}[ht]
\centering
\includegraphics[width=0.85\textwidth]{figures/seam_23_57_87_soft.pdf}
\caption{Shared-maturity comparison for two independently calibrated relaxed blocks on the 2026-04-21 surfaces.
The terminal SPX representation of the first block and the initial SPX representation of the second block differ by up to approximately $1.1$ volatility points and by $0.3$ volatility points on average over the invertible strike range at the displayed matched operating point.
An exact stitched model with a prescribed common marginal has no such duplication; the figure diagnoses the separate numerical selection made by the relaxed block solves.}
\label{fig:market_seam}
\end{figure}

\begin{table}[ht]
\centering
\caption{Conditional post-projection ladder for the two relaxed blocks.
The middle columns give the reported blockwise dispersion diagnostics and the last two columns give the maximum and mean shared-smile differences over invertible strikes, in volatility points.
The values describe a finite post-projection path, not exact stitched feasible laws.}
\label{tab:seam_tightness}
\small
\begin{tabular}{@{}lcccc@{}}
\toprule
Conditional target & $\edisp^{A}$ & $\edisp^{B}$ & seam max (vp) & seam mean (vp) \\
\midrule
None & $1.3$ & $7.6$ & $0.91$ & $0.29$ \\
$10^{-1}$ & $1.3$ & $3.0$ & $1.22$ & $0.35$ \\
$10^{-2}$ & $3.1\times10^{-1}$ & $2.3\times10^{-1}$ & $1.13$ & $0.34$ \\
$10^{-3}$ & $3.2\times10^{-2}$ & $2.0\times10^{-2}$ & $3.62$ & $1.54$ \\
$10^{-4}$ & $3.8\times10^{-3}$ & $4.4\times10^{-3}$ & $5.40$ & $2.55$ \\
Exact conditional post-projection & $6.5\times10^{-14}$ & $9.3\times10^{-14}$ & $5.53$ & $2.62$ \\
\bottomrule
\end{tabular}
\end{table}

\subsection{Finite-Budget Penalty Path}\label{subsec:market_penalty}

\Cref{tab:market_penalty} reports the shared-penalty sweep on the five-surface numerical grid.
Across the tested path the worst implied-volatility error remains below $0.70$ volatility points.
Between $\lambda=1$ and the representative operating point $\lambda=10^4$, the thresholded bulk dispersion statistic falls from $5.39$ to $6.3\times10^{-2}$, approximately an $86$-fold reduction, while the bulk martingale statistic falls from $7.3\times10^{-3}$ to $7.6\times10^{-4}$.
The sequence is not monotone at every tested value, and the $\lambda=3\times10^4$ row is worse than the $\lambda=10^4$ row on both conditional summaries.
We therefore interpret $\lambda=10^4$ as a useful finite-budget operating point, not as an asymptotic optimum.

\begin{table}[ht]
\centering
\caption{Finite-budget conditional-penalty path on the smoothed 2026-04-21 five-surface inputs and the $1.2\times10^6$-cell grid.
The middle column is the worst implied-volatility error over the five fitted surfaces, in volatility points; $\edisp$ and $\emart$ are the thresholded bulk statistics of~\eqref{eq:massw}.
Each row is an independent fixed-budget solve.}
\label{tab:market_penalty}
\small
\begin{tabular}{@{}rccc@{}}
\toprule
$\lambda$ & worst smile (vp) & $\edisp$ & $\emart$ \\
\midrule
$1$        & $0.54$ & $5.39$ & $7.3\times10^{-3}$ \\
$10$       & $0.61$ & $2.31$ & $3.1\times10^{-3}$ \\
$100$      & $0.58$ & $1.14$ & $1.7\times10^{-3}$ \\
$300$      & $0.61$ & $0.37$ & $1.0\times10^{-3}$ \\
$1{,}000$  & $0.61$ & $8.9\times10^{-2}$ & $6.0\times10^{-4}$ \\
$3{,}000$  & $0.54$ & $7.1\times10^{-2}$ & $7.9\times10^{-4}$ \\
$10{,}000$ & $0.55$ & $6.3\times10^{-2}$ & $7.6\times10^{-4}$ \\
$30{,}000$ & $0.58$ & $8.7\times10^{-2}$ & $8.7\times10^{-4}$ \\
\bottomrule
\end{tabular}
\end{table}

At $\lambda=10^4$, the $10\%$-of-peak filter retains $78.1\%$ of the conditioning mass in the first transition and $44.3\%$ in the second.
When every positive-mass cell is included instead, the corresponding worst-over-transition mass-weighted statistics are $1.94$ for dispersion and $5.0\times10^{-3}$ for martingality, compared with the displayed bulk values $6.3\times10^{-2}$ and $7.6\times10^{-4}$.
This tail sensitivity is why we use the qualified term \emph{thresholded bulk residual} and do not treat the table as an exact-feasibility test.

\begin{figure}[ht]
\centering
\includegraphics[width=0.9\textwidth]{figures/p2_sweetspot_paper.pdf}
\caption{Finite-budget penalty path on the 2026-04-21 five-surface inputs.
The plot displays the worst fitted-smile error and the thresholded bulk conditional statistics against $\lambda$.
The selected $\lambda=10^4$ point balances the reported diagnostics within the tested schedule; local reversals preclude interpreting the curve as a monotone convergence path.}
\label{fig:market_penalty}
\end{figure}

\begin{figure}[ht]
\centering
\includegraphics[width=\textwidth]{figures/p2_sweetspot_smile_25f6de8.pdf}
\caption{Five fitted implied-volatility smiles at the representative $\lambda=10^4$ operating point.
The worst reported error over the selected strikes is $0.55$ volatility points, with thresholded bulk statistics $(\edisp,\emart)=(6.3\times10^{-2},7.6\times10^{-4})$.
The statistic describes the aligned five-axis approximation and the retained conditioning cells specified in \cref{subsec:market_setup}.}
\label{fig:market_smiles}
\end{figure}

Separating the two conditional penalties clarifies which family drives this calculation.
With the dispersion penalty pinned at $100$, the martingale statistic changes only from about $1.7\times10^{-3}$ to $1.5\times10^{-3}$ as $\lambda_{\mathrm{mart}}$ rises to $10^4$.
With the martingale penalty pinned at $100$, the dispersion statistic falls from $5.14$ to $5.7\times10^{-2}$.
Thus dispersion is the binding soft family for this implementation and dataset; this is an empirical attribution, not a universal property of SPX--VIX calibration.

\begin{table}[ht]
\centering
\caption{Per-family penalty ablation on the same numerical problem.
One conditional penalty is varied while the other is fixed at $100$.
The residual is the varied family's thresholded bulk statistic.}
\label{tab:market_family}
\small
\begin{tabular}{@{}rcc@{}}
\toprule
sweep $\lambda_{\mathrm{mart}}$ (dispersion pinned) & $\emart$ & worst smile (vp) \\
\midrule
$1$--$1{,}000$ & $\sim1.7\times10^{-3}$ & $0.58$ \\
$10{,}000$ & $1.5\times10^{-3}$ & $0.58$ \\
\midrule
sweep $\lambda_{\mathrm{disp}}$ (martingale pinned) & $\edisp$ & worst smile (vp) \\
\midrule
$1$ & $5.14$ & $0.54$ \\
$100$ & $1.14$ & $0.58$ \\
$1{,}000$ & $8.9\times10^{-2}$ & $0.61$ \\
$10{,}000$ & $5.7\times10^{-2}$ & $0.61$ \\
\bottomrule
\end{tabular}
\end{table}

\subsection{Conditional Post-Projection Diagnostics}\label{subsec:market_projection}

A separate diagnostic on the same $1.2\times10^6$-cell grid starts from the law obtained without a conditional penalty and applies a conditional-only Newton projection.
\Cref{tab:market_projection} compares the three independently recorded regimes using the common maximum diagnostics displayed in the corresponding heatmaps.
The conditional-only projection reduces the displayed conditional errors below $10^{-6}$ but increases the worst displayed smile error to $10.8$ volatility points.
This is a conditional projection applied after the no-conditional-penalty solve, not evidence that the exact joint affine system is infeasible.

\begin{table}[ht]
\centering
\caption{Common heatmap diagnostics for the no-conditional-penalty, soft-$\lambda=100$, and conditional-only projection regimes.
The conditional columns are maximum relative errors over the populated cells displayed by that experiment; they are different from the mass-weighted summaries in \cref{tab:market_penalty}.}
\label{tab:market_projection}
\small
\begin{tabular}{@{}lccc@{}}
\toprule
Regime & worst smile (vp) & dispersion max & martingale max \\
\midrule
No conditional penalty & $0.531$ & $19.52$ & $6.06\times10^{-2}$ \\
Soft conditional penalty, $\lambda=100$ & $0.531$ & $10.32$ & $1.52\times10^{-2}$ \\
Conditional-only projection & $10.786$ & $1.4\times10^{-7}$ & $2.9\times10^{-10}$ \\
\bottomrule
\end{tabular}
\end{table}

\begin{figure}[ht]
\centering
\includegraphics[width=\textwidth]{figures/p2_heatmaps.pdf}
\caption{Spatial diagnostics for the no-conditional-penalty, soft-$\lambda=100$, and conditional-only projection regimes.
The conditional projection makes the displayed conditional maps nearly exact while concentrating the deterioration in fitted-smile regions; the soft run leaves the quote fit near its initial level and reports the remaining conditional discrepancy.}
\label{fig:market_heatmaps}
\end{figure}

The recorded alternating-projection trace gives a complementary finite-budget diagnostic.
Within each exact round, the conditional update reduces its own residual and the subsequent quote-row update reopens it.
The post-round endpoints nevertheless trend downward, so the trace does not prove a limit cycle or asymptotic nonconvergence.
Over the recorded budgets, the best joint residual was $1.51\times10^{-3}$ for exact alternation and $1.68\times10^{-4}$ for the augmented run, approximately a nine-fold difference.
The augmented trace also contains local reversals; the comparison is therefore a finite-budget result only.

\begin{figure}[ht]
\centering
\includegraphics[width=0.49\textwidth]{figures/p2_oscillation_exact.pdf}\hfill
\includegraphics[width=0.49\textwidth]{figures/p2_oscillation_augmented.pdf}
\caption{Recorded finite-budget traces for exact family alternation (left) and the augmented method (right).
The sawtooth behavior shows the handoff between quote and conditional updates.
The reported comparison concerns the best joint residual reached within the tested budgets and is not an asymptotic convergence claim.}
\label{fig:market_oscillation}
\end{figure}

\subsection{Misaligned-Date Calibration on a Merged Timeline}\label{subsec:market_misaligned}

A separate experiment used a 2026-05-13 snapshot with three SPX surfaces and two VIX surfaces.
The seven-axis merged-timeline model includes intermediate SPX states at the two VIX dates.
The saved seven-axis calculation has $1.12\times10^6$ cells and reports a maximum conditional-row residual of $9.4\times10^{-4}$ and a maximum $A_h$-row residual of $1.0\times10^{-4}$.
Compared with the five-axis interpolated baseline in \cref{tab:market_misaligned}, the seven-axis merged-timeline model improves the reported maximum and mean implied-volatility errors on that instance.

The two 30-day VIX spans overlap by two days.
Consequently this calculation lies outside the non-overlap assumption of \cref{prop:merged_reference}; it is evidence that the enlarged numerical state space can be useful, but it does not validate the exact reference construction proved there.
It also uses a different snapshot from the primary 2026-04-21 experiments.

\begin{table}[ht]
\centering
\caption{Five-axis and seven-axis merged-timeline fits on the 2026-05-13 misaligned-date instance.
The seven-axis model includes intermediate SPX states at the two VIX dates.
Entries are maximum and mean implied-volatility errors over the selected strikes, in volatility points.
Because the two VIX spans overlap, the calculation is not covered by \cref{prop:merged_reference}.}
\label{tab:market_misaligned}
\small
\begin{tabular}{@{}lcccccr@{}}
\toprule
 & SPX$_{T_1}$ & VIX$_1$ & SPX$_{T_2}$ & VIX$_2$ & SPX$_{T_3}$ & cells \\
\midrule
$5$D interpolated, max & $0.221$ & $1.228$ & $0.119$ & $0.472$ & $0.069$ & $1.20\times10^{6}$ \\
$5$D interpolated, mean & $0.103$ & $0.261$ & $0.053$ & $0.118$ & $0.025$ & $1.20\times10^{6}$ \\
\midrule
$7$D merged timeline, max & $0.059$ & $0.711$ & $0.071$ & $0.240$ & $0.027$ & $1.12\times10^{6}$ \\
$7$D merged timeline, mean & $0.009$ & $0.135$ & $0.012$ & $0.077$ & $0.008$ & $1.12\times10^{6}$ \\
\bottomrule
\end{tabular}
\end{table}

\subsection{Empirical Scaling}\label{subsec:market_timing}

Finally, a fixed $720$-sweep schedule was run on single-threaded grids ranging from approximately $78{,}000$ to $1.2$ million cells.
Over this tested range, a log--log fit gave runtime proportional to $N^{1.04}$ with $R^2=0.9998$.
Recorded wall time ranged from approximately $80$ seconds to $22$ minutes, while problem construction on the largest grid took less than $6$ seconds.
This empirical exponent is consistent with the linear-in-$d$ per-sweep accounting of \cref{app:complexity}; it is not an iteration-complexity theorem and no claim is made beyond the tested hardware and schedule.

\begin{figure}[ht]
\centering
\includegraphics[width=0.55\textwidth]{figures/timing_scaling_5d.pdf}
\caption{Recorded single-thread wall time against five-axis grid size for a fixed $720$-sweep schedule.
The fitted relation over the tested range is proportional to $N^{1.04}$ with $R^2=0.9998$; the largest $1.2$-million-cell run took approximately $22$ minutes.}
\label{fig:market_timing}
\end{figure}

\section{Conclusion and Future Work}\label{sec:conclusion}

We developed a global state-space formulation for multi-maturity SPX--VIX calibration.
For its finite-dimensional implementation, we proposed an augmented-Bregman mirror-descent scheme.
The exact theory draws a sharp line between feasibility and identification: local and global nonemptiness are equivalent, but stitching is a block-preserving Markovian projection that can discard dependence relevant to multi-period pricing and risk.
With the standard Markov reference, entropy selects this stitched law as the minimum-information completion; global coupling becomes economically operative when cross-period information or objectives are supplied.

The numerical contribution addresses a different issue.
On finite supports, marginal and conditional rows can be incompatible at the requested tolerances.
The proposed scheme preserves the fit to observable quote moments and exposes the remaining martingale and dispersion residuals along a finite-budget penalty path.
The controlled synthetic experiment shows why this explicit allocation is preferable to hiding an unavoidable discrepancy in prescribed marginals.
The market-data calculations then illustrate the same allocation on smoothed SPX and VIX surfaces: fitted-smile errors remain below $0.70$ volatility points across the reported penalty sweep while the thresholded bulk conditional diagnostics improve substantially.
The independent-block seam and conditional post-projection experiments are practical finite-budget diagnostics, not evidence against exact local--global feasibility equivalence.

\paragraph{Future directions.} A central challenge is that the full coupling tensor grows exponentially in the number of maturities $m$, while the theoretical advantage of the global model derives precisely from conditioning on the full history.
Promising directions that preserve the non-Markovian character include:
\emph{truncated-history conditioning} (windowed constraints on the recent path, exponential in the window but independent of $m$);
\emph{parametric history dependence} (conditioning on low-dimensional path summary statistics, e.g.\ a weighted average of past volatilities, optionally pre-encoded into the reference measure);
\emph{continuous dual parameterization} (optimizing marginal potentials and conditional multipliers as functions via stochastic gradient ascent on the dual, avoiding materialization of the tensor); and GPU/sparse-grid acceleration.
Further directions include robust price bounds over compatible block gluings, history-dependent reference measures, cross-period calibration targets, a rigorous treatment of overlapping VIX windows, risk sensitivities, and extensions to other asset classes with analogous identification gaps.

\clearpage
\appendix
\section*{Appendix}

\section{Proofs of the Structural and Variational Results}\label[appendix]{app:proofs}

\subsection{Band compatibility}\label[appendix]{app:proof_band_compatibility}

\begin{proof}[Proof of \cref{prop:band_compatibility}]
If $\mu\in\cP_{\mathrm{full}}^{\mathrm{band}}$, set $\alpha_i=\Law_\mu(S_i)$ and $\beta_i=\Law_\mu(V_i)$.
The adjacent triple law $\Law_\mu(S_i,V_i,S_{i+1})$ is locally feasible by the tower property, so $(\alpha_i,\beta_i,\alpha_{i+1})\in\mathfrak R_i$ for every $i$.
The sequence is common across all blocks because each $S_i$ is one random variable under $\mu$.

Conversely, suppose one sequence $(\alpha_1,\beta_1,\ldots,\beta_{m-1},\alpha_m)$ belongs to all adjacent relations.
Choose $\nu_i\in\cP(\alpha_i,\beta_i,\alpha_{i+1})$ for each $i$ and disintegrate each later block with respect to its first SPX coordinate.
The product
\[
\nu_1(ds_1,dv_1,ds_2)
\prod_{i=2}^{m-1}\nu_i(dv_i,ds_{i+1}\mid s_i)
\]
is well-defined because adjacent blocks share $\alpha_i$.
It preserves every selected marginal and, by the same conditional-expectation argument as in \cref{app:proof_feasibility_equivalence}, satisfies (C2)--(C3).
It therefore belongs to $\cP_{\mathrm{full}}^{\mathrm{band}}$.
Separate nonemptiness of the $\mathfrak R_i$ need not yield one common intermediate $\alpha_i$, which proves the final statement.
\end{proof}

\subsection{Feasibility equivalence and Markovization}\label[appendix]{app:proof_feasibility_equivalence}

\begin{proof}[Proof of \cref{thm:feasibility_equivalence}]
Let $\mu\in\cP_{\mathrm{full}}$.
Since $\sigma(S_i,V_i)\subseteq\cF_i$, the tower property gives
\begin{align*}
\Embb^\mu[S_{i+1}\mid S_i,V_i]
&=\Embb^\mu[\Embb^\mu[S_{i+1}\mid\cF_i]\mid S_i,V_i]=S_i,\\
\Embb^\mu[L(S_{i+1}/S_i)\mid S_i,V_i]
&=\Embb^\mu[\Embb^\mu[L(S_{i+1}/S_i)\mid\cF_i]\mid S_i,V_i]=V_i^2.
\end{align*}
Thus $\nu_i^\mu:=\Law_\mu(S_i,V_i,S_{i+1})$ belongs to $\cP_i$ for every $i$.

The product in~\eqref{eq:markovization} preserves $\nu_1^\mu$.
Inductively, if it has the same $S_i$ marginal as $\mu$, adjoining the kernel $\mu(dv_i,ds_{i+1}\mid s_i)$ gives the same $(S_i,V_i,S_{i+1})$ law $\nu_i^\mu$; hence it also preserves the next $S_{i+1}$ marginal.
Every adjacent triple law is therefore preserved.
Moreover, the product factorization gives
\[
\Law_{\mathsf M\mu}(S_{i+1}\mid\cF_i)
=\Law_{\nu_i^\mu}(S_{i+1}\mid S_i,V_i),
\]
so the two local conditional identities imply (C2)--(C3) under $\mathsf M\mu$.
Thus $\mathsf M\mu\in\cP_{\mathrm{stitch}}\subseteq\cP_{\mathrm{full}}$.

Finally, if every $\cP_i$ is nonempty, choose $\nu_i\in\cP_i$ and glue them by the product in~\eqref{eq:stitched_class}; their prescribed shared SPX marginals agree, and the preceding argument proves that the resulting law is globally feasible.
The reverse implications follow from the block preservation just established.
\end{proof}

\subsection{Strict inclusion and non-identification}\label[appendix]{app:proof_strict_inclusion}

\begin{proof}[Proof of \cref{thm:strict_inclusion}]
The inclusion follows from \cref{thm:feasibility_equivalence}.
For strictness, take $m=3$, set $S_0=S_1=100$, and let $V_1$ be equally likely to equal $v_L=0.15$ or $v_H=0.55$.
For $(v,a)=(v_L,8)$ or $(v_H,30)$, set
\[
q(v,a):=\frac{v^2}{-(2/\tau)\log(1-(a/100)^2)}\in(0,1/2)
\]
and, conditional on $V_1=v$, assign probabilities $q(v,a)$, $1-2q(v,a)$, and $q(v,a)$ to $S_2=100-a$, $100$, and $100+a$.
Symmetry gives $\Embb[S_2\mid S_1,V_1]=S_1$, while
\[
q(v,a)\left[-\frac{2}{\tau}\log(1-a/100)-\frac{2}{\tau}\log(1+a/100)\right]=v^2,
\]
which gives the first dispersion identity.

Set $V_2=0.20$ when $V_1=v_L$ and $V_2=0.60$ when $V_1=v_H$ and, conditionally on the history, let
\[
S_3=S_2(1+\varepsilon\delta(V_2)),
\qquad
\delta(v):=\sqrt{1-e^{-v^2\tau}},
\]
where $\varepsilon$ is conditionally uniform on $\{-1,1\}$.
Then $\Embb[S_3\mid\cF_2]=S_2$ and
\[
\Embb[L(S_3/S_2)\mid\cF_2]
=-\frac{1}{\tau}\log(1-\delta(V_2)^2)=V_2^2,
\]
so the resulting finite-support law $\mu$ belongs to $\cP_{\mathrm{full}}$.

Both $V_1$ branches reach $S_2=100$ with positive probability, but $V_2$ identifies the branch there.
Hence $V_2\not\perp V_1\mid S_2$, so~\eqref{eq:stitched_ci} fails and $\mu\notin\cP_{\mathrm{stitch}}$.
Its Markovization belongs to $\cP_{\mathrm{stitch}}$ and preserves both adjacent triple laws by \cref{thm:feasibility_equivalence}.
Because the two finite laws differ, the indicator of any atom on which their masses differ is a bounded separating payoff.
Finally, equality $\mu=\mathsf M\mu$ is equivalent to equality of the full-history kernels and the $S_i$-conditional kernels in~\eqref{eq:markovization}, which is exactly~\eqref{eq:stitched_ci}.
For $m>3$, append degenerate feasible periods.
\end{proof}

\subsection{Information and entropy identities}\label[appendix]{app:proof_information_loss}

\begin{proof}[Proof of \cref{prop:information_loss}]
Disintegrate $\mu$ chronologically as
\[
\mu(d\mathbf s,d\mathbf v)
=\nu_1^\mu(ds_1,dv_1,ds_2)
\prod_{i=2}^{m-1}\mu(dv_i,ds_{i+1}\mid H_{i-1},s_i),
\]
whereas~\eqref{eq:markovization} replaces each displayed continuation kernel by $\mu(dv_i,ds_{i+1}\mid s_i)$.
The chain rule for relative entropy therefore gives
\begin{align*}
D_{\mathrm{KL}}(\mu\Vert\mathsf M\mu)
&=\sum_{i=2}^{m-1}
\Embb^\mu\!\left[
\log\frac{d\mu((V_i,S_{i+1})\mid H_{i-1},S_i)}
{d\mu((V_i,S_{i+1})\mid S_i)}
\right]\\
&=\sum_{i=2}^{m-1}I_\mu((V_i,S_{i+1});H_{i-1}\mid S_i).
\end{align*}
Each summand is nonnegative and vanishes exactly when the corresponding conditional-independence relation~\eqref{eq:stitched_ci} holds.
\end{proof}

\subsection{Duality and existence}\label[appendix]{app:proof_duality}

\begin{proof}[Proof of \cref{thm:duality}]
If $\mu\in\cP_{\mathrm{full}}$, then every bounded predictable gain in~\eqref{eq:portfolio_value} has zero expectation by (C2)--(C3), while (C1) prices the static terms at their initial costs.
Thus $\Embb^\mu[\Pi]=\mathrm{Cost}(\Pi)$ for every $\Pi\in\cU_{\mathrm{full}}$, ruling out a pointwise nonnegative portfolio of negative cost.

Conversely, for a law satisfying (C1), the identity
\[
\Embb^\mu[\Delta_{S,i}(\mathbf S^i,\mathbf V^i)(S_{i+1}-S_i)]=0
\]
for every bounded $\cF_i$-measurable $\Delta_{S,i}$ is equivalent to (C2); the analogous identity with $L(S_{i+1}/S_i)-V_i^2$ is equivalent to (C3).
The dynamic multipliers therefore dualize exactly the two conditional families.
On a finite state space, infeasibility of the resulting linear system is equivalent by Farkas' lemma to the existence of multipliers whose combined payoff is nonnegative on every state and whose initial cost is negative.
This is the finite-dimensional form of the martingale-transport separation argument; continuum analogues require the additional hypotheses discussed in \cref{subsec:duality}, as in \citet{Beiglbck2013,Beiglbck2017} and the SPX--VIX formulation of \citet{guyon2020joint}.
\end{proof}

\subsection{Entropic existence and the reference measure}\label[appendix]{app:proof_entropic_af}

\begin{proof}[Proof of \cref{thm:entropic_af}]
Take a minimizing sequence in the asserted weakly compact entropy sublevel.
It has a weakly convergent subsequence; weak closedness of $\cP_{\mathrm{full}}$ keeps the limit feasible, and lower semicontinuity of relative entropy makes that limit a minimizer.
Relative entropy is strictly convex on laws dominated by $\bar\mu$, while $\cP_{\mathrm{full}}$ is convex, so two distinct finite-entropy minimizers cannot exist.
The resulting feasible law rules out global arbitrage by \cref{thm:duality}.
On a finite allowed support, the simplex is compact, the affine constraint set is closed, and a strictly positive reference dominates every feasible vector; hence nonemptiness supplies all of the stated conditions.
\end{proof}

\begin{proof}[Proof of \cref{prop:ref_props}]
The iterated product in~\eqref{eq:complete_ref} is a probability measure by construction.
The lognormal kernel $T_i$ satisfies
\[
\Embb[S_{i+1}\mid S_i=s_i,V_i=v_i]=s_i
\]
and, since $\log(S_{i+1}/S_i)=v_i\sqrt\tau G-v_i^2\tau/2$,
\[
\Embb[L(S_{i+1}/S_i)\mid S_i=s_i,V_i=v_i]=v_i^2.
\]
The product construction makes these identities valid conditional on the full history and draws each $V_i$ from $\mu_{V_i}$ independently of the preceding state; it also starts from $\mu_{S_1}$.
Its support is the product-kernel support generated by these factors, not necessarily the entire ambient space.
On a finite allowed grid, strictly positive reference weights dominate every law on that grid.
In the continuum, neither topological full support nor equality of supports implies absolute continuity, which is why \cref{thm:entropic_af} assumes a feasible finite-entropy law separately.
\end{proof}

\subsection{Missing intermediate SPX marginals}\label[appendix]{app:proof_stitch_illposed}

\begin{proof}[Proof of \cref{prop:stitch_illposed}]
A standard three-variable block beginning at $T_i^V$ takes the initial SPX marginal as part of its input.
When $T_i^V\notin\mathcal T_S$, the supplied SPX option surfaces specify no such marginal.
Consequently the block is not determined by the stated market inputs: one must either select a marginal by an additional interpolation rule or optimize over it as an endogenous variable in an enlarged state space.
\end{proof}

\subsection{Merged-timeline reference}\label[appendix]{app:proof_merged_reference}

\begin{proof}[Proof of \cref{prop:merged_reference}]
Starting from the prescribed law of $S_{t_1}$, successively adjoining the stated VIX draws and SPX transition kernels defines a probability law by finite iterated disintegration.
Its $V_i$ marginals are the prescribed ones by construction.
Partition a non-overlapping VIX span into merged-timeline subintervals of lengths $\Delta t_k$ and let $d_{i,k}$ be the corresponding forward ratios.
Conditional on $\cF_{T_i^V}$ and $V_i$, define independent increments by
\[
\frac{S_{t_{k+1}}}{d_{i,k}S_{t_k}}
=\exp\!\left(V_i\sqrt{\Delta t_k}\,G_k-\tfrac12V_i^2\Delta t_k\right),
\qquad G_k\sim\mathcal N(0,1).
\]
Each increment has conditional mean one, so their product is a forward-adjusted martingale over the span.
Since $\sum_k\Delta t_k=\tau_i$ and the Gaussian terms have conditional mean zero,
\[
\Embb\!\left[-\frac{2}{\tau_i}\sum_k
\log\!\left(\frac{S_{t_{k+1}}}{d_{i,k}S_{t_k}}\right)
\middle|\cF_{T_i^V}\right]
=\frac{1}{\tau_i}\sum_k V_i^2\Delta t_k
=V_i^2.
\]
The sum of logs is the log of the forward-adjusted endpoint ratio.
Non-overlap ensures that no subinterval is assigned two different VIX levels, while the stipulated kernels on uncovered intervals preserve the martingale property.
The complete chronological product therefore has all the properties claimed.
\end{proof}

\subsection{KL selection of the stitched completion}\label[appendix]{app:proof_kl_markovization}

\begin{proof}[Proof of \cref{prop:kl_markovization}]
Apply the relative-entropy chain rule to the full-history kernels of $\mu$ and the $S_i$-conditional kernels of $\bar\mu$.
For each $i\ge2$, insert the intermediate kernel $\mu(dv_i,ds_{i+1}\mid s_i)$ into the log-density ratio.
The terms comparing the full-history kernel with this intermediate kernel sum to $D_{\mathrm{KL}}(\mu\Vert\mathsf M\mu)$ by \cref{prop:information_loss}.
The remaining terms depend only on the adjacent block laws; because $\mu$ and $\mathsf M\mu$ share those laws, together with the first block, they sum to $D_{\mathrm{KL}}(\mathsf M\mu\Vert\bar\mu)$.
This proves~\eqref{eq:kl_markovization}.
Since $\mathsf M\mu$ is feasible by \cref{thm:feasibility_equivalence}, Markovization weakly decreases the objective, strictly so whenever $D_{\mathrm{KL}}(\mu\Vert\mathsf M\mu)>0$.
Therefore every finite-entropy minimizer is stitched.
\end{proof}

\subsection{Finite-grid penalty limit}\label[appendix]{app:proof_penalty_limit}

\begin{proof}[Proof of \cref{prop:penalty_limit}]
The set $\mathcal H$ is a closed subset of the finite-dimensional simplex and is therefore compact.
The objective in~\eqref{eq:penalty_problem} is continuous on $\mathcal H$ and strictly convex because $\bar\pi$ is positive, proving existence and uniqueness of $\pi_\lambda$.
For any $\pi_\star\in\argmin_{\mathcal H}r$, optimality gives
\[
D_{\mathrm{KL}}(\pi_\lambda\Vert\bar\pi)+\lambda r(\pi_\lambda)
\leq
D_{\mathrm{KL}}(\pi_\star\Vert\bar\pi)+\lambda r_\star.
\]
Relative entropy is nonnegative, while it is bounded above on the finite simplex when $\bar\pi$ is positive.
Consequently $0\le r(\pi_\lambda)-r_\star\le D_{\mathrm{KL}}(\pi_\star\Vert\bar\pi)/\lambda$, so $r(\pi_\lambda)\to r_\star$.
Compactness gives limit points, all in $\argmin_{\mathcal H}r$ by continuity.
Taking $\pi_\star$ to be the minimum-KL element of that set and using the same optimality inequality shows that every limit point has no larger relative entropy.
Strict convexity on the convex set $\argmin_{\mathcal H}r$ makes this element unique, so the entire sequence converges to it.
The two stated cases follow from whether $r_\star$ is zero.
\end{proof}

\section{The Online-Mirror-Descent View and Infeasibility}\label[appendix]{app:omd_theory}

This appendix collects the mirror-descent interpretation of the solver, its relation to exact row projection, and a controlled analysis of what each scheme sacrifices under infeasibility.

\paragraph{Mirror descent and exact row projection.} Mirror descent is the standard first-order scheme that moves along the (sub)gradient of the objective in the geometry induced by a strictly convex mirror map and projects back with the associated Bregman divergence~\citep{nemirovski1983problem,beck2003mirror}.
With the negative-entropy mirror map on the simplex its update is the multiplicative (exponentiated-gradient) reweighting $\pi\odot\exp(-\alpha g)$ followed by renormalization, and the exact KL/Bregman projection onto a single linear constraint is the Sinkhorn-type diagonal rescaling used for the hard families~\citep{Benamou2015Iterative}.
The augmented-Bregman solver of \cref{alg:augmented} uses the residual gradient~\eqref{eq:aug_gradient} for the conditional block and cyclic Bregman--Newton corrections for the $A_h$ rows.

\begin{proposition}[Exact row projection as an entropic mirror step]\label{prop:sinkhorn_omd}
Let $p$ be strictly positive on a finite support and suppose $\{q\in\Delta_d:a^\top q=b\}$ meets the relative interior of the simplex.
The unique solution of
\[
\min_{q\in\Delta_d}\left\{D_{\mathrm{KL}}(q\Vert p):a^\top q=b\right\}
\]
has the form
\[
q_j(\theta)=\frac{p_j e^{\theta a_j}}{\sum_k p_k e^{\theta a_k}},
\qquad a^\top q(\theta)=b.
\]
Thus replacing the first-order update~\eqref{eq:mirror_step} by this solved KL proximal step gives the usual exponential Bregman projection~\citep{csiszar1975,bauschke1997}.
Visiting all rows cyclically with these solved proximal steps gives the classical cyclic Bregman/Sinkhorn iteration.
For the separately defined pure-penalty problems~\eqref{eq:penalty_problem}, the exact-constraint limit is given by \cref{prop:penalty_limit}.
\end{proposition}

\begin{proof}
With multipliers $\lambda_0$ and $\lambda_1$ for normalization and the affine row, stationarity of
\[
\sum_j q_j\log(q_j/p_j)+\lambda_0\left(\sum_jq_j-1\right)
+\lambda_1(a^\top q-b)
\]
gives $q_j=p_j e^{\theta a_j}/Z(\theta)$ after absorbing constants and writing $\theta=-\lambda_1$.
The derivative of $\log Z(\theta)$ is $a^\top q(\theta)$ and its second derivative is $\operatorname{Var}_{q(\theta)}(a)$.
Hence the multiplier is unique when $a$ is nonconstant and $b$ lies in the relative interior of its attainable interval; the constant-row case is immediate.
Cyclically visiting these solved subproblems yields the exact Bregman iteration.
\end{proof}

\paragraph{What each scheme sacrifices under infeasibility.} The hard/soft split of \cref{subsec:discrete_formulation} is justified by a controlled synthetic experiment with an exact, tunable feasibility knob.
We discretize two centered Gaussian laws on the common $41$-point grid specified in \cref{sec:experiments} and impose both marginals together with the martingale rows on the $25$ source states above the stated mass threshold.
The variance ordering required by Strassen's theorem~\citep{strassen1965existence} motivates the two cases, and a direct LP verifies that the variance-expanded affine system is feasible while the variance-contracted system is infeasible even with the tail rows omitted.
Within each instance, three schemes differing in their priority allocation (cyclic row projection, an all-soft OMD penalty, and a marginal-priority hybrid) are compared in \cref{tab:feasibility}.
All runs start from the product of the two prescribed discrete marginals.
Rows are stored and visited in this fixed order: all source-marginal rows, all target-marginal rows, then the active martingale rows.
For a hard row, Newton iterations solve its scalar exponential tilt to absolute residual $10^{-13}$, subject to a cumulative log-tilt cap of $8$; normalization is restored after every row.
For soft rows, one batch exponentiated-gradient update with penalty $200$, base step $0.5$, and log-tilt cap $1$ follows the hard-row sweep.
The marginal-priority hybrid uses hard marginal rows and a soft martingale block without dual ascent; it is a controlled priority-allocation test, not an execution of \cref{alg:augmented}.
When feasible, cyclic row projection ends with marginal and conditional residuals $5.1\times10^{-6}$ and $8.8\times10^{-14}$, while the finite-budget hybrid holds the marginal residual to $2.8\times10^{-5}$.
When \emph{infeasible}, the schemes differ in what they sacrifice: cyclic row projection nearly satisfies the conditional rows but leaves marginal residual $1.9\times10^{-2}$, whereas the hybrid keeps the marginals about $25$ times tighter ($7.4\times10^{-4}$) and exposes a conditional residual of $1.1\times10^{-3}$.
The numerical ratio is specific to this order and finite budget; the qualitative design lesson is to state explicitly which family receives priority.

\begin{table}[ht]
\centering
\caption{Synthetic feasibility knob (discretized Gaussian marginals, $n=41$).
Entries are median maximum absolute affine residuals over the final $10\%$ of each run; cyclic row projection and the marginal-priority hybrid use $800$ sweeps, while all-soft OMD uses $6000$.
Under the stated row order and schedule, the marginal-priority method keeps the prescribed marginals about $25\times$ tighter than cyclic row projection.}
\label{tab:feasibility}
\small
\begin{tabular}{@{}lcccc@{}}
\toprule
 & \multicolumn{2}{c}{\textbf{Feasible} ($\rho_{\mathrm{var}}\approx1.3$)} & \multicolumn{2}{c}{\textbf{Infeasible} ($\rho_{\mathrm{var}}\approx0.7$)} \\
\cmidrule(lr){2-3}\cmidrule(lr){4-5}
\textbf{Scheme} & marginal & conditional & marginal & conditional \\
\midrule
Cyclic row projection & $5.1\times10^{-6}$ & $8.8\times10^{-14}$ & $1.9\times10^{-2}$ & $1.3\times10^{-5}$ \\
OMD (all soft)         & $3.1\times10^{-2}$ & $1.2\times10^{-3}$  & $4.5\times10^{-2}$ & $1.2\times10^{-3}$ \\
Marginal-priority      & $2.8\times10^{-5}$ & $2.5\times10^{-4}$  & $\mathbf{7.4\times10^{-4}}$ & $1.1\times10^{-3}$ \\
\bottomrule
\end{tabular}
\end{table}

\section{Complexity of the Augmented-Bregman Solver}\label[appendix]{app:complexity}

We collect the cost accounting for \cref{alg:augmented}.
Throughout, $d$ is the number of cells in the innovation-coordinate tensor $\pi$; for common grid sizes in the aligned discretization, $d=n_Sn_V^{m-1}n_Z^{m-1}$.
We count arithmetic operations per sweep.
The estimates are informal but capture the dependence on tensor size and maturity count.

\paragraph{The operator is sparse.} The conditional rows are structurally sparse.
Fix a transition $i$; the martingale (C2) and dispersion (C3) constraints partition the cells of $\pi$ into fibers indexed by the history $(s_1,v_1,\ldots,s_i,v_i)$, and each cell belongs to exactly one such fiber for each of the two families.
Each active cell enters at most one martingale and one dispersion fiber per transition, so $A_c$ has at most $2(m{-}1)d$ nonzeros.
The penalty gradient $A_c^\top(y+\lambda(A_c\pi-b_c))$, two sparse matrix--vector products and an element-wise reweighting, therefore costs $O(md)$.
Each normalization, forward, or selected option-payoff row scans at most $d$ cells; if $R_h$ such rows are visited, their cyclic corrections cost $O(R_hd)$ per sweep.

\paragraph{Cost per sweep and in total.}
A single sweep costs $O((m+R_h)d)$, which is linear in the tensor size for a fixed maturity count and fixed number of quote rows.
With the fixed $K_{\mathrm{out}}\times K_{\mathrm{in}}$ schedule of \cref{alg:augmented}, the measured work is therefore
$O(K_{\mathrm{out}}K_{\mathrm{in}}(m+R_h)d)$.
We do not infer a global iteration-complexity bound for the complete augmented scheme from the standard mirror-descent bound for a fixed convex objective, because the iteration includes cyclic row corrections, dual ascent, clipping, and an optional changing penalty.

\paragraph{Comparison.} The schemes differ in their per-sweep cost and finite-budget residual behavior on a coupled problem.
Stitching decomposes into $m{-}1$ independent three-variable solves and is cheaper, at the price of the conditional-independence assumption identified in \cref{thm:strict_inclusion}.
A conditional Newton projection adds a small solve per active fiber, while the augmented method uses sparse matrix--vector operations and row corrections.
The exponential growth of $d$ in $m$ is intrinsic to the full-history coupling and is what the scaling directions of \cref{sec:conclusion}, such as truncated-history conditioning and continuous dual parameterization, are meant to mitigate.

\section*{Acknowledgments}

The authors thank Junhyung Lyle Kim, Rudy Raymond, Ruslan Shaydulin and Rob Otter for valuable feedback and discussions.
We also acknowledge our colleagues at the Global Technology Applied Research Center of JPMorganChase for support throughout this work.

\section*{Disclaimer}
This paper was prepared for informational purposes by the Global Technology Applied Research center of JPMorgan Chase \& Co. This paper is not a product of the Research Department of JPMorgan Chase \& Co. or its affiliates. Neither JPMorgan Chase \& Co. nor any of its affiliates makes any explicit or implied representation or warranty and none of them accept any liability in connection with this paper, including, without limitation, with respect to the completeness, accuracy, or reliability of the information contained herein and the potential legal, compliance, tax, or accounting effects thereof. This document is not intended as investment research or investment advice, or as a recommendation, offer, or solicitation for the purchase or sale of any security, financial instrument, financial product or service, or to be used in any way for evaluating the merits of participating in any transaction.

\bibliographystyle{plainnat}
\bibliography{refs}

\end{document}
